% Canon X-07 -- Référence technique
% Synthèse de la 2e partie (« Mystères du X-07 », chap. 7 à 10, p. 65-116)
% du livre « Les mystères du Canon X-07 » (scan x_07_mysteres_x-07.pdf).
% Compilation : pdflatex X07_mysteres_reference.tex  (deux passes pour la table des matières)

\documentclass[a4paper,10pt]{article}

\usepackage[utf8]{inputenc}
\usepackage[T1]{fontenc}
\usepackage{lmodern}
\usepackage[french]{babel}
\usepackage[margin=2cm]{geometry}
\usepackage{array}
\usepackage{booktabs}
\usepackage{longtable}
\usepackage{enumitem}
\usepackage[table]{xcolor}
\usepackage[hidelinks]{hyperref}

\setlength{\parindent}{0pt}
\setlength{\parskip}{4pt}
\setlist{nosep}

\newcommand{\hex}[1]{\texttt{#1}}
\newcommand{\reg}[1]{\textsf{#1}}
\newcommand{\pg}[1]{\texorpdfstring{\textcolor{gray}{\small(p.~#1)}}{}}
% Remarque du rédacteur : coquille ou incohérence relevée dans le livre
\newcommand{\nota}[1]{\textcolor{red!60!black}{\small\textbf{N.d.R.}~#1}}

% Environnement pour les fiches de routines ROM
\newenvironment{routines}
  {\small\begin{longtable}{@{}>{\ttfamily\bfseries}p{1.1cm}>{\raggedright}p{4.0cm}>{\raggedright}p{5.0cm}>{\raggedright}p{3.1cm}>{\raggedright\arraybackslash}p{2.1cm}@{}}
   \toprule
   \normalfont\bfseries Adr. & \textbf{Entrée} & \textbf{Action} & \textbf{Sortie} & \textbf{Modifiés}\\
   \midrule\endhead
   \bottomrule\endfoot}
  {\end{longtable}}

% Fiche d'identité d'un programme
\newcommand{\fiche}[3]{%
  \begin{center}\small
  \begin{tabular}{@{}ll@{}}
  \toprule
  Implantation & #1 \\
  Longueur & #2 \\
  Lancement & #3 \\
  \bottomrule
  \end{tabular}
  \end{center}}

% Listing désassemblé : adresse, octets, étiquette, instruction, commentaire
\newenvironment{asm}
  {\small\begin{longtable}{@{}>{\ttfamily}l>{\ttfamily}l>{\ttfamily}l>{\ttfamily}l>{\raggedright\arraybackslash}p{6.2cm}@{}}
   \toprule
   \normalfont\bfseries Adr. & \normalfont\bfseries Octets & \normalfont\bfseries Étiq. & \normalfont\bfseries Instruction & \textbf{Commentaire}\\
   \midrule\endhead
   \bottomrule\endfoot}
  {\end{longtable}}

\title{\textbf{Canon X-07 -- Référence technique}\\[4pt]
\large Synthèse des 2\textsuperscript{e} et 3\textsuperscript{e} parties\\
« Mystères du X-07 » et « Applications »\\
du livre \emph{Les mystères du Canon X-07}}
\author{}
\date{}

\begin{document}
\maketitle

\begin{center}
\begin{minipage}{0.85\textwidth}\small
Ce document reprend les pages 65 à 135 du livre. La première moitié (chapitres 7
à 10) rassemble sous forme de tableaux les informations techniques :
architecture, carte mémoire, ports d'entrée/sortie, sous-processeur T6834,
codage du BASIC, zone système, routines de la ROM et périphériques. La
seconde reprend les huit programmes en langage machine de la
3\textsuperscript{e} partie « Applications », avec leurs listings vérifiés
par désassemblage. Les numéros de page \pg{xx} renvoient au livre. Sauf mention
contraire, les valeurs sont en hexadécimal. Les coquilles et incohérences
relevées dans le livre sont signalées par \nota{}. Les notes « vérifié dans
la ROM » s'appuient sur le désassemblage du vidage de la ROM BASIC
(\texttt{ROM\_x07/x07.bin}, implanté en \hex{B000}--\hex{FFFF}).
\end{minipage}
\end{center}

\tableofcontents
\newpage

%======================================================================
\part*{2\textsuperscript{e} partie du livre : Mystères du X-07 \pg{65-116}}
\addcontentsline{toc}{part}{\texorpdfstring{2\textsuperscript{e}}{2e} partie du livre : Mystères du X-07}

\section{Architecture générale \pg{65-67}}

Le X-07 est entièrement en technologie CMOS et fonctionne sur 4 piles de 1,5\,V.
Il utilise deux processeurs :
\begin{itemize}
  \item un \textbf{NSC800} (compatible Z80), le processeur principal ;
  \item un sous-processeur \textbf{Toshiba T6834}, qui gère le clavier,
        l'affichage, les caractères graphiques, l'horloge, l'alarme, le buzzer,
        l'alimentation et le contrôle des piles.
\end{itemize}
Le NSC800 envoie ses commandes au T6834 par le circuit d'E/S \textbf{HD61L202F}.
Une fois que le T6834 a accusé réception, le processeur principal continue.
La mémoire de l'écran LCD n'est pas accessible directement par le Z80 : elle
est adressée par le T6834.

\subsection{Composants (fig.~18)}
\begin{center}\small
\begin{tabular}{@{}ll@{}}
\toprule
Fonction & Composant \\
\midrule
Processeur principal & NSC800 \\
RAM interne 8\,Ko & 4 $\times$ HM6117FLP \\
Support d'extension RAM & 8\,Ko \\
Connecteur de carte (« utility card ») & jusqu'à 16\,Ko \\
ROM masquée & 16\,Ko + 4\,Ko (BASIC, 20\,Ko au total) \\
Circuit d'E/S & HD61L202F \\
Sous-processeur & T6834, avec 2\,Ko de RAM (HM6117FLP) \\
Écran & LCD 20 colonnes $\times$ 4 lignes \\
Pilotes d'affichage & 2 $\times$ T6781 \\
Bus d'extension & vers l'interface vidéo X-720 \\
Ports & parallèle (Centronics), série, cassette, buzzer, coupleur optique \\
\bottomrule
\end{tabular}
\end{center}

Extensions possibles : 16\,Ko de RAM (support ou carte), carte ROM Canon
de 8\,Ko, ROM de 4\,Ko et RAM vidéo de l'interface X-720.

%======================================================================
\section{Carte mémoire \pg{67-69}}

L'espace d'adressage est de 64\,Ko. Au démarrage, une ROM d'amorçage occupe
l'adresse \hex{0000}. Elle est ensuite masquée par la RAM interne.
L'adressage du support RAM et de la carte s'inverse lorsqu'une RAM
additionnelle de 8\,Ko est présente, pour que l'espace reste continu quand la
carte est retirée.

\begin{center}\small
\begin{tabular}{@{}lll@{}}
\toprule
Adresses & Taille & Contenu \\
\midrule
\hex{0000}--\hex{1FFF} & 8\,Ko & RAM interne (ROM d'amorçage superposée au démarrage) \\
\hex{2000}--\hex{3FFF} & 8\,Ko & RAM de la carte (si une ROM est connectée) \emph{ou} support RAM (si une RAM est connectée) \\
\hex{4000}--\hex{5FFF} & 8\,Ko & support ROM \emph{ou} RAM de la carte \\
\hex{6000}--\hex{7FFF} & 8\,Ko & ROM de la carte \\
\hex{8000}--\hex{97FF} & 6\,Ko & RAM vidéo (X-720) \\
\hex{A000}--\hex{AFFF} & 4\,Ko & ROM TV (X-720) \\
\hex{B000}--\hex{FFFF} & 20\,Ko & ROM BASIC \\
\bottomrule
\end{tabular}
\end{center}

\subsection{Détection des cartes ROM \pg{67-69}}
La carte est détectée par logiciel. À l'initialisation, la ROM cherche le
mot-clé \texttt{love} entre \hex{2000} et \hex{B000}, par pas de \hex{800}
(2048 octets). Les 6 octets qui suivent le mot-clé forment 3 adresses de
sous-programmes :
\begin{itemize}
  \item les deux premières sont toujours exécutées ;
  \item la troisième est exécutée après un \texttt{SLEEP}.
\end{itemize}

%======================================================================
\section{Ports d'entrée/sortie du HD61L202F \pg{69-74}}\label{sec:ports}

Le circuit d'E/S gère le port Centronics, le port série (UART), le buzzer,
la cassette, le générateur de bauds, le coupleur optique et le dialogue avec
le T6834. Il occupe les ports \hex{F0} à \hex{F7}.

\subsection{Ports en écriture (fig.~21)}
\begin{center}\small
\begin{tabular}{@{}l>{\raggedright}p{3.2cm}*{8}{c}@{}}
\toprule
Port & Rôle & D7 & D6 & D5 & D4 & D3 & D2 & D1 & D0 \\
\midrule
\hex{F0} & contrôle des interruptions & OFFRQ & SETPS & BIE & ALMIE & XTIE & XRIE & TIE & RIE \\
\hex{F1} & XBTR : donnée vers le T6834 & \multicolumn{8}{c}{octet} \\
\hex{F2} & bauds, poids faible & \multicolumn{8}{c}{bits 0--7 du compteur} \\
\hex{F3} & bauds, poids fort & -- & -- & -- & -- & \multicolumn{4}{c}{bits 8--11} \\
\hex{F4} & modes & SETBC & BRGST & CNTR & LEO & MD1 & MD0 & BZON & REM \\
\hex{F5} & RAZ des interruptions & RALRM & CSTB & CSFT & RBGI & RTDRE & RRDRF & RXTRE & RXRRF \\
\hex{F6} & configuration UART & -- & IR & RTS & ER & SBRK & RXE & -- & TXEN \\
\hex{F6}' & format série (2\textsuperscript{e} écriture) & S2 & -- & EP & PEN & -- & L1 & B2 & FX04 \\
\hex{F7} & donnée série à émettre & \multicolumn{8}{c}{octet} \\
\bottomrule
\end{tabular}
\end{center}

\paragraph{Port \hex{F0}.} OFFRQ : coupure de l'alimentation ; SETPS : mode
économie (éteint) ; BIE : interruption du générateur de bauds ; ALMIE :
interruption d'alarme ; XTIE / XRIE : interruptions du bus vers le T6834 ;
TIE / RIE : interruptions d'émission / de réception série.
Attention : modifier \hex{F0} sans précaution peut bloquer la machine.

\paragraph{Ports \hex{F2}/\hex{F3}.} Compteur de 12 bits du générateur de
bauds : les 8 bits de poids faible dans \hex{F2}, les 4 bits de poids fort
dans les bits bas de \hex{F3}.

\paragraph{Port \hex{F4}.} SETBC\,=\,1 charge \hex{F2}/\hex{F3} dans le
générateur ; BRGST\,=\,1 lance le comptage, 0 l'arrête ; CNTR : bit envoyé
vers la prise parallèle ; LEO\,=\,1 envoie 38,4\,kHz au coupleur optique ;
BZON : buzzer ; REM\,=\,1 active la télécommande du magnétophone.

\begin{center}\small
\begin{tabular}{@{}ccl@{}}
\toprule
MD1 & MD0 & Mode (fig.~22) \\
\midrule
0 & 0 & RAZ du compteur de bauds, arrêt \\
0 & 1 & liaison série \\
1 & 0 & cassette (signal sur la prise K7) \\
1 & 1 & son sur le buzzer \\
\bottomrule
\end{tabular}
\end{center}

\paragraph{Port \hex{F6} (écriture).} IR : initialisation de l'UART ;
RTS\,=\,1 met la broche RTS du HD61L202F à 0 ; ER : remise à zéro des
indicateurs d'erreur PE, OE et FE ; SBRK : force la ligne à 0 (break) ;
RXE : autorise la réception ; TXEN : autorise l'émission.

\subsection{Ports en lecture (fig.~23)}
\begin{center}\small
\begin{tabular}{@{}l>{\raggedright}p{3.2cm}*{8}{c}@{}}
\toprule
Port & Rôle & D7 & D6 & D5 & D4 & D3 & D2 & D1 & D0 \\
\midrule
\hex{F0} & interruptions & TRQ1 & TRQ0 & BIE & ALMIE & XTIE & XRIE & TIE & RIE \\
\hex{F1} & XBRR : donnée du T6834 & \multicolumn{8}{c}{octet} \\
\hex{F2} & indicateurs & CBSY & ALRM & AUX1 & AUX2 & 0 & BGI & XTRE & XRRF \\
\hex{F4} & modes & FCAS & BRGST & CNTR & LEO & MD1 & MD0 & BZON & REM \\
\hex{F6} & état de l'UART & CTS & -- & FE & OE & PE & TXE & RXRDY & TXRDY \\
\hex{F7} & donnée série reçue & \multicolumn{8}{c}{octet} \\
\bottomrule
\end{tabular}
\end{center}
Les ports \hex{F3} et \hex{F5} ne sont pas utilisés en lecture.

\begin{itemize}
  \item \hex{F0} : TRQ0\,=\,1 si la touche BREAK est enfoncée ; TRQ1\,=\,1 est
        positionné par le T6834 : le NSC800 ne peut alors pas passer en mode
        sauvegarde.
  \item \hex{F2} : CBSY : port parallèle occupé ; ALRM\,=\,1 : piles
        épuisées (testé toutes les 4 minutes par le sous-processeur) ;
        AUX2 : pile de la carte mémoire ; AUX1\,=\,0 : interrupteur de la prise
        d'extension sur RAM ; BGI : sortie du générateur de bauds ;
        XTRE\,=\,1 quand le T6834 a lu XBTR ; XRRF\,=\,1 quand le T6834 a
        écrit XBRR.
  \item \hex{F4} : FCAS : reflet du signal lu sur la cassette.
  \item \hex{F6} : la donnée série est émise si CTS\,=\,0 et TXEN\,=\,1 ;
        FE : bit de stop manquant ; OE : donnée reçue avant la lecture de la
        précédente (overrun) ; PE : erreur de parité ; TXE : registre
        d'émission vide.
\end{itemize}
\nota{Le livre donne « T*RDY : réception prête, R*RDY : transmission prête » ;
les libellés semblent inversés (RXRDY = réception prête, TXRDY = émission prête).}

%======================================================================
\section{Le sous-processeur T6834 \pg{74-86}}

Le T6834 accepte 70 commandes (codes \hex{01} à \hex{46}), envoyées par le
NSC800 à travers le HD61L202F. Une commande peut être suivie de paramètres.
Si nécessaire, le T6834 répond par une suite d'octets.

\subsection{Liste des commandes (§ 7.6.1 et fig.~24)}
« Env. » est le nombre d'octets envoyés, code de commande compris ; « Rép. »
est le nombre d'octets de la réponse ; $n$ désigne une chaîne de longueur
variable, terminée par \hex{00}.

{\small
\begin{longtable}{@{}>{\ttfamily}c>{\raggedright}p{3.6cm}cc>{\raggedright\arraybackslash}p{7.4cm}@{}}
\toprule
\normalfont Code & Nom & Env. & Rép. & Description \\
\midrule\endhead
\bottomrule\endfoot
01 & TIME CALL & 1 & 8 & lit l'heure : année (2 octets), mois, date, jour de la semaine, heure, minutes, secondes \\
02 & STICK & 1 & 1 & touche curseur : \hex{30} aucune, \hex{31} haut, \hex{32} droite, \hex{36} bas, \hex{37} gauche \\
03 & STRIG & 1 & 1 & 0 si F6 est enfoncée, sinon \hex{FF} \\
04 & STRIG1 & 1 & 1 & 0 si la barre d'espace est enfoncée, sinon \hex{FF} \\
05 & RAM READ & 3 & 1 & lit la RAM du T6834 (adresse : octet haut, octet bas) \\
06 & RAM WRITE & 4 & 0 & écrit dans la RAM du T6834 (adresse haute, basse, donnée) \\
07 & SCROLL SET & 3 & 0 & première et dernière ligne de défilement \\
08 & SCROLL EXEC & 1 & 0 & fait défiler l'écran (le curseur ne bouge pas) \\
09 & LINE CLEAR & 2 & 0 & efface la ligne indiquée \\
0A & TIME SET & 9 & 0 & règle l'horloge : seconde, minute, heure, jour, date, mois, année (2 octets) \\
0B & calcul du jour & 1 & 0 & calcule et règle le jour de la semaine d'après TIME SET \\
0C & ALARM SET & 9 & 0 & règle l'alarme, même format que \hex{0A} (le 2\textsuperscript{e} paramètre est ignoré) \\
0D & BUZZER OFF & 1 & 0 & pas de buzzer à l'heure d'alarme \\
0E & BUZZER ON & 1 & 0 & buzzer à l'heure d'alarme \\
0F & BIT IMAGE & 2 & 120 & renvoie une ligne de l'écran (0 à 3) au format « bit image » \\
10 & POINT & 3 & 1 & 0 si le point $(X,Y)$ est allumé, sinon \hex{FF} \\
11 & PSET & 3 & 0 & allume le point $(X,Y)$ \\
12 & PRESET & 3 & 0 & éteint le point $(X,Y)$ \\
13 & PEOR & 3 & 0 & inverse le point $(X,Y)$ \\
14 & LINE & 5 & 0 & trace un segment ; paramètres $X_1,Y_1,X_2,Y_2$ \\
15 & CIRCLE & 4 & 0 & trace un cercle ; paramètres $X,Y,r$ \\
16 & UDK WRITE & 2+$n$ & 0 & définit une touche de fonction (numéro 1 à 12, chaîne : 3 caractères de repérage + 38 caractères ; 42 caractères pour F6 et F12) \\
17 & UDK READ & 2 & $n$ & lit la définition d'une touche de fonction \\
18 & UDK ON & 1 & 0 & le T6834 commence à envoyer les UDK \\
19 & UDK OFF & 1 & 0 & le T6834 cesse d'envoyer les UDK \\
1A & UDC WRITE & 10 & 0 & définit un caractère (codes \hex{80}--\hex{9F} et \hex{E0}--\hex{FF}) : code + 8 octets (6$\times$8 points) \\
1B & UDC READ & 2 & 8 & lit un motif de caractère (codes \hex{20}--\hex{FF}) \\
1C & UDC INIT & 1 & 0 & rétablit les caractères par défaut \\
1D & SP WRITE & 1+$n$ & 0 & écrit le programme de lancement (511 octets au plus, terminé par \hex{00}) \\
1E & SP WRITE CONT & 1+$n$ & 0 & ajoute à la suite de \hex{1D} \\
1F & SP ON & 1 & 0 & autorise le programme de lancement \\
20 & SP OFF & 1 & 0 & invalide le programme de lancement \\
21 & SP READ & 1 & $n$ & lit le programme de lancement \\
22 & ON STATE & 1 & 1 & bit 0\,=\,1 : OFF/ON ; bit 6\,=\,1 : SLEEP ON \\
23 & OFF REQUEST & 1 & 0 & coupe l'alimentation \\
24 & LOCATE & 4 & 0 & place le curseur en $(X,Y)$ et affiche le caractère (\hex{00} : aucun caractère) \\
25 & CURSOR ON & 1 & 0 & affiche le curseur \\
26 & CURSOR OFF & 1 & 0 & cache le curseur \\
27 & TEST KEY & 3 & 1 & teste une ligne de la matrice du clavier (2 octets, voir fig.~25) \\
28 & TEST CHR & 2 & 1 & 0 si la touche du caractère ASCII est enfoncée, sinon \hex{FF} (sans SHIFT) \\
29 & RAZ secondes & 1 & 0 & remet les secondes à zéro \\
2A & DATE SET & 2 & 0 & écrit la date \\
2B & DISPLAY OFF & 1 & 0 & coupe le signal commun du LCD (la mémoire d'affichage reste accessible) \\
2C & DISPLAY ON & 1 & 0 & rétablit l'affichage \\
2D & KEY BUFFER CLEAR & 1 & 0 & vide le tampon du clavier (sans toucher aux UDK ni au SP) \\
2E & CLS & 1 & 0 & efface l'écran \\
2F & HOME & 1 & 0 & curseur en $(0,0)$ \\
30 & UDK DISP ON & 1 & 0 & affiche les étiquettes des UDK sur la ligne 3 \\
31 & UDK DISP OFF & 1 & 0 & efface les étiquettes des UDK \\
32 & REPEAT KEY ON & 1 & 0 & active la répétition des touches \\
33 & REPEAT KEY OFF & 1 & 0 & désactive la répétition \\
34 & KANA SET & 1 & 0 & assigne les caractères KANA aux UDK \\
35 & UDK CONT WRITE & 2+$n$ & 0 & ajoute des caractères à la définition d'une touche de fonction \\
36 & ALARM READ & 1 & 8 & lit l'alarme (même format que \hex{01}) \\
37 & BUZZER ZERO & 1 & 1 & 0 si le compteur est à 0, sinon \hex{FF} \\
38 & CLICK OFF & 1 & 0 & désactive le « clic » des touches \\
39 & CLICK ON & 1 & 0 & active le « clic » des touches \\
3A & LOCATE CLOSE & 1 & 0 & ramène le curseur dans la zone de défilement (après \hex{24}) \\
3B & KEY OUT ON & 1 & 0 & autorise l'entrée des touches vers l'unité centrale \\
3C & KEY OUT OFF & 1 & 0 & bloque l'entrée des touches vers l'unité centrale \\
3D & SP AUTO ON & 1 & 0 & exécute le programme de lancement à la mise sous tension \\
3E & SP AUTO OFF & 1 & 0 & n'exécute plus le programme de lancement à la mise sous tension \\
3F & REQUEST SLEEP & 1 & 0 & coupe l'alimentation en mode SLEEP \\
40 & UDK INIT & 1 & 0 & rétablit les touches de fonction par défaut \\
41 & CHAR WRITE & 9 & 0 & affiche un motif 6$\times$8 au curseur (8 octets, format \hex{1A}) ; le curseur avance \\
42 & CHAR READ & 1 & 8 & renvoie les 8 octets du motif sous le curseur \\
43 & SCANR & 3 & 2 & points de la zone 6$\times$8 au curseur (format \hex{1B}) \\
44 & SCANL & 3 & 2 & nombre de points allumés à gauche de $(X-1,Y)$ \\
45 & TIME CHK & 1 & 1 & vérifie l'horloge : bit 0 secondes, 1 minutes, 2 heures, 3 jour, 4 mois, 5 année (1901--2099) ; un bit à 1 signale une erreur ; bits 6 et 7 à 0 \\
46 & ALM CHK & 1 & 1 & vérifie l'alarme \\
\end{longtable}
}
\nota{Le livre écrit « 38 CLICK OFF déclenche » et « 39 CLICK ON enclenche » ;
on retient OFF = désactivé, ON = activé.}

\nota{Vérifié dans la ROM : \texttt{STICK} (\hex{D7E3}) envoie
\hex{82} (commande 02, bit 7 = aucun paramètre) et garde la réponse
\texttt{AND \$0F}. \texttt{STRIG(0)} (\hex{D7F8}) envoie \hex{84}
(barre d'espace) et \texttt{STRIG(1)} \hex{83} (F6) ; la réponse est
inversée par \texttt{CPL} : $-1$ si la touche est enfoncée, 0 sinon. Les deux
fonctions vident ensuite les tampons (\hex{C0BD}).}

\subsection{Matrice du clavier (fig.~25)}
Pour la commande \hex{27} (TEST KEY), on envoie la valeur de la ligne
d'échantillonnage sur 2 octets ; le T6834 renvoie l'état des colonnes
K1 à K8.

\begin{center}\small
\begin{tabular}{@{}ll*{8}{c}@{}}
\toprule
Ligne & Valeur & K1 & K2 & K3 & K4 & K5 & K6 & K7 & K8 \\
\midrule
A1 & \hex{0200} & SFT & CTR & & & & & & \\
BZ & \hex{0100} & Home & & GRP & Num & OFF & & & ON \\
S8 & \hex{0080} & \multicolumn{8}{l}{L, +, *, \}, ), \textbar, =, -- (ponctuation)} \\
S7 & \hex{0040} & \multicolumn{8}{l}{>, ?/$\times$, Yen, CR, O, P, @, \{ (ponctuation)} \\
S6 & \hex{0020} & F1 & F2 & F3 & F4 & F5 & F6 & & \\
S5 & \hex{0010} & 1 & 2 & 3 & 4 & 5 & 6 & 7 & 8 \\
S4 & \hex{0008} & Q & W & E & R & T & Y & U & I \\
S3 & \hex{0004} & A & S & D & F & G & H & J & K \\
S2 & \hex{0002} & Z & X & C & V & B & N & M & \\
S1 & \hex{0001} & INS & DEL & $\rightarrow$ & $\leftarrow$ & $\uparrow$ & $\downarrow$ & SPC & \\
\bottomrule
\end{tabular}
\end{center}
Les lignes S7 et S8 sont difficiles à lire sur le scan : on n'en donne que
les touches principales.

\subsection{RAM du T6834 (§ 7.6.2)}
Le T6834 dispose de 2\,Ko de RAM, qu'on lit et écrit par les commandes
\hex{05}/\hex{06} ou par les routines \hex{E348}/\hex{E334}.

\begin{center}\small
\begin{tabular}{@{}l>{\raggedright\arraybackslash}p{11.5cm}@{}}
\toprule
Adresse & Contenu \\
\midrule
\hex{C006} & bit 0 : ONSTATE ; bit 6 : état SLEEP ; bit 7 : mode ALARM simulé \\
\hex{C00C} & mode du clavier : 0 normal, 1 KANA, 2 graphique, 3 numérique \\
\hex{C01A} & bit 0\,=\,1 : alarme validée \\
\hex{C040}--\hex{C04E} & heure courante, une valeur tous les 2 octets : secondes, minutes, heures, jour de la semaine, date, mois, année \\
\hex{C050}--\hex{C05E} & alarme, même format \\
\hex{C080} & indicateur du curseur \\
\hex{C082} & indicateur des touches de fonction utilisateur \\
\hex{C084} & compteur du buzzer \\
\hex{C092} & compteur de la minuterie d'alimentation \\
\hex{C096} & fin de la zone de défilement \\
\hex{C09E} & début de la zone de défilement \\
\hex{C800}--\hex{C9FF} & touches de fonction (512 octets, voir ci-dessous) \\
\hex{CA00}--\hex{CBFF} & caractères définis : codes \hex{80}--\hex{9F} en $\hex{CA00}+(c-\hex{80})\times 8$ ; codes \hex{E0}--\hex{FF} en $\hex{CB00}+(c-\hex{E0})\times 8$ \\
\hex{CC00}--\hex{CCFF} & tampon du clavier (256 octets, 127 touches) \\
\hex{CE00}--\hex{CFFF} & programme de lancement (512 octets, \texttt{START\$}) \\
\hex{E000}--\hex{E1FE} & mémoire de l'écran LCD (511 octets) \\
\bottomrule
\end{tabular}
\end{center}

\paragraph{Adresses des touches de fonction (fig.~26).}
\begin{center}\small
\begin{tabular}{@{}*{6}{lc}@{}}
\toprule
F1 & \hex{C800} & F2 & \hex{C82A} & F3 & \hex{C85A} & F4 & \hex{C87E} & F5 & \hex{C8A8} & F6 & \hex{C8D2} \\
F7 & \hex{C900} & F8 & \hex{C92A} & F9 & \hex{C954} & F10 & \hex{C97E} & F11 & \hex{C9A8} & F12 & \hex{C9D2} \\
\bottomrule
\end{tabular}
\end{center}
\nota{Avec un pas de \hex{2A}, on attendrait \hex{C854} pour F3 (comme F9 = \hex{C954}) ;
la valeur \hex{C85A} du livre est probablement une coquille.}

\subsection{Utilisation du T6834 en langage machine (§ 7.6.3)}

La ROM fournit deux routines principales :
\begin{itemize}
  \item \hex{E428} : envoie une commande \textbf{sans paramètre ni réponse}.
        Charger le code dans \reg{A}, puis \texttt{CALL \$E428}. Modifie
        \reg{A}, \reg{F} et \reg{C} (\reg{L} selon le §\,7.6.3).
  \item \hex{C92F} : envoie une commande \textbf{avec paramètres et réponse}.
        \reg{A} = code de la commande ; \reg{HL} = adresse des paramètres ;
        \reg{B} = nombre de paramètres ; \reg{C} = nombre d'octets de la
        réponse ; \reg{DE} = tampon de réponse (en général \hex{026E}).
        Modifie \reg{AF}, \reg{BC}, \reg{DE} et \reg{HL}. Si $B=0$, le bit 7
        de \reg{A} doit être positionné.
\end{itemize}

\paragraph{Exemple : faire clignoter l'écran.}
\begin{verbatim}
        LD   B,$FF
loop:   LD   A,$2B       ; DISPLAY OFF
        CALL $E428
        LD   A,$2C       ; DISPLAY ON
        CALL $E428
        DJNZ loop
        RET
\end{verbatim}
\nota{Le livre écrit \texttt{LD B,\$2B} à la 2\textsuperscript{e} ligne ; il faut \texttt{LD A,\$2B}.}

\paragraph{Routine d'envoi d'un octet (livre).} Cette routine envoie au T6834
l'octet contenu dans \reg{E}, sans modifier aucun registre. La routine ROM
\hex{C9C0} attend que le T6834 soit prêt.
\begin{verbatim}
        PUSH AF
        PUSH DE
        PUSH BC
        LD   C,$F1       ; port de sortie vers le T6834
        CALL $C9C0       ; attente : T6834 prêt ?
        LD   A,($026C)   ; image du port F0 (INTIMAG)
        OR   $80
        OUT  ($F0),A
        OUT  (C),E       ; octet -> port F1
        LD   A,$02
        OUT  ($F5),A
        POP  BC
        POP  DE
        POP  AF
        RET
\end{verbatim}
\nota{Le listing du livre porte \texttt{CALL \$C0C9} et \texttt{POP HL} ; le
code hexadécimal du programme de démonstration (\texttt{CD C0 C9}, puis \texttt{F1})
montre qu'il s'agit de \texttt{CALL \$C9C0} et de \texttt{POP AF}.}

\paragraph{Programme de démonstration (BASIC).} Il implante la routine
ci-dessus en \hex{1C25}, efface l'écran (\texttt{CALL \$CE9E}), trace un cercle
(commande \hex{15}, $X=\hex{3C}$, $Y=\hex{0F}$, $r=\hex{0A}$), puis fait
clignoter l'écran 255 fois (commandes \hex{2B}/\hex{2C}).
\begin{verbatim}
 5 REM programme BASIC implantant une routine en LANGAGE MACHINE
10 CLS : PRINT "Je charge les codes ... Un instant, SVP !"
20 FOR X = &H1C00 TO &H1C3E
30 READ A$ : POKE X,VAL("&H" + A$)
40 NEXT X
45 REM donnees de la routine LANGAGE MACHINE
50 DATA CD,9E,CE,1E,15,CD,25,1C
60 DATA 1E,3C,CD,25,1C,1E,0F,CD
70 DATA 25,1C,1E,0A,CD,25,1C,06
80 DATA FF,1E,2B,CD,25,1C,1E,2C
90 DATA CD,25,1C,10,F4,F5,C5,D5
92 DATA 0E,F1,CD,C0,C9,3A,6C,02
95 DATA F6,80,D3,F0,ED,59,3E,02
98 DATA D3,F5,D1,C1,F1,C9,00,00
\end{verbatim}
Lancement : \texttt{EXEC \&H1C00}.

%======================================================================
\section{Le BASIC en mémoire \pg{87-90}}

Les programmes BASIC commencent à l'adresse \textbf{1363} (décimal), soit
\hex{0553}. La zone de 0 à 1362 est la zone système (§\,\ref{sec:zone}).
Chaque mot-clé est codé par un octet entre \hex{80} et \hex{FF} (« token »).

\paragraph{Format d'une ligne.}
\begin{center}\small
\begin{tabular}{@{}ll@{}}
\toprule
Octets & Contenu \\
\midrule
2 & adresse de la ligne suivante (poids faible, poids fort) \\
2 & numéro de ligne (poids faible, poids fort) \\
$n$ & tokens et caractères ASCII \\
1 & \hex{00} : fin de ligne \\
\bottomrule
\end{tabular}
\end{center}

\paragraph{Exemple du livre.} Le programme \texttt{1 REM **}, \texttt{200 CLS:A\$="C7"},
\texttt{9999 PRINT"END","C7":END} donne, à partir de l'adresse 1363 :
\texttt{92 05 01 00 8E 20 2A 2A 00 | 6A 05 C8 00 A6 3A 41 24 DD 22 43 37 22 00 | 7C 05 0F 27 9F \dots{} 80 00}
(en hexadécimal ; le livre donne ces valeurs en décimal).

\subsection{Table des mots-clés et points d'entrée en ROM (fig.~27 et 28)}
L'interpréteur BASIC occupe \hex{B000}--\hex{FFFF}. Le tableau donne le token
de chaque mot-clé et l'adresse de la routine correspondante en ROM.
Les adresses \hex{0087}--\hex{009C} sont des « hooks » en RAM (zone système),
redirigés par l'interface X-720.

{\small
\setlength{\tabcolsep}{4pt}
\begin{longtable}{@{}lcc|lcc|lcc@{}}
\toprule
Mot-clé & Tok. & Entrée & Mot-clé & Tok. & Entrée & Mot-clé & Tok. & Entrée \\
\midrule\endhead
\bottomrule\endfoot
END & 80 & D2B0 & FOR & 81 & F3F7 & NEXT & 82 & D30C \\
DATA & 83 & F66A & INPUT & 84 & F82F & DIM & 85 & B005 \\
READ & 86 & F86B & LET & 87 & F685 & GOTO & 88 & F621 \\
RUN & 89 & F5B7 & IF & 8A & F797 & RESTORE & 8B & D296 \\
GOSUB & 8C & F610 & RETURN & 8D & F63D & REM & 8E & F666 \\
STOP & 8F & D2C0 & ELSE & 90 & F666 & TR & 91 & D1B7 \\
MOTOR & 92 & E0B9 & DEFSTR & 93 & F54E & DEFINT & 94 & F551 \\
DEFSNG & 95 & F554 & DEFDBL & 96 & F557 & LINE & 97 & F7C6 \\
ERROR & 98 & F78C & RESUME & 99 & F734 & OUT & 9A & E98C \\
ON & 9B & F6F1 & LPRINT & 9C & E9DE & DEFFN & 9D & FC96 \\
POKE & 9E & FC80 & PRINT & 9F & EA05 & CONT & A0 & D178 \\
LIST & A1 & FE75 & LLIST & A2 & FE6B & CLEAR & A3 & D1D0 \\
CIRCLE & A4 & 009C & CONSOLE & A5 & E243 & CLS & A6 & CE9E \\
COLOR & A7 & 008A & EXEC & A8 & CEB8 & LOCATE & A9 & E212 \\
PSET & AA & 008D & PRESET & AB & 0090 & OFF & AC & C5B1 \\
SLEEP & AD & C59B & DIR & AE & E5E3 & DELETE & AF & E682 \\
FSET & B0 & BF56 & PAINT & B1 & 0099 & LOAD & B2 & DF38 \\
SAVE & B3 & DEFE & INIT & B4 & E699 & ERASE & B5 & E3CF \\
BEEP & B6 & C2CD & CLOAD & B7 & DF30 & CSAVE & B8 & DEF9 \\
NEW & B9 & D214 & FN & BC & FCB5 & ERL & BE & FAC4 \\
ERR & BF & FAB6 & STRING\$ & C0 & D835 & INSTR & C1 & DCD6 \\
INKEY\$ & C2 & C103 & INP & C3 & E931 & VARPTR & C4 & E9A7 \\
USR & C5 & E9BB & SNS & C6 & E961 & ALM\$ & C7 & D928 \\
DATE\$ & C8 & D8CE & TIME\$ & C9 & D868 & START\$ & CA & DA99 \\
FONT\$ & CB & D739 & KEY\$ & CC & DA39 & SCREEN & CD & 0087 \\
NOT & CF & FC1A & STEP & D0 & F537 & EQV & D9 & FC72 \\
AND & D6 & FC68 & SGN & DF & C9F1 & INT & E0 & CC23 \\
ABS & E1 & C9DC & FRE & E2 & DDD7 & POS & E3 & FC8C \\
SQR & E4 & B694 & RND & E5 & B771 & LOG & E6 & B609 \\
EXP & E7 & B6DD & COS & E8 & B530 & SIN & E9 & B549 \\
TAN & EA & B593 & ATN & EB & B5AC & PEEK & EC & FC7A \\
CINT & ED & CAE0 & CSNG & EE & CB08 & CDBL & EF & CB90 \\
FIX & F0 & CC14 & LEN & F1 & D72D & HEX\$ & F2 & D538 \\
STR\$ & F3 & D53D & VAL & F4 & DCB3 & ASC & F5 & D81A \\
CHR\$ & F6 & D828 & TKEY & F7 & D7CC & LEFT\$ & F8 & DC5A \\
RIGHT\$ & F9 & DC8A & MID\$ & FA & DC93 & CSRLIN & FB & CEAF \\
STICK & FC & D7E3 & STRIG & FD & D7F8 & POINT & FE & 0093 \\
\end{longtable}
\nota{Le livre donne \hex{BF36} pour FSET. La table de répartition de la ROM
(\hex{F0D6}) donne \hex{BF56} (\hex{BF36} est au milieu de la conversion des
nombres) ; ses voisines INIT, ERASE, LOAD, SAVE concordent avec le livre.
Vérifié dans la ROM et en appelant le gestionnaire depuis du code compilé.}
}
Codes de contrôle de la fig.~27 : BELL (sonnerie), BS (retour arrière), LF
(saut de ligne), LU (ligne précédente), CLS, CR, flèches de déplacement du
curseur, DC1/DC2 (passage texte/graphique sur la X-710), SI/SO (caractères
de la X-710 en ROM externe/interne).

%======================================================================
\section{Zone de travail du système \pg{90-95}}\label{sec:zone}

La zone système occupe \hex{0000}--\hex{0552} (1363 octets) :
\begin{itemize}
  \item \hex{0000}--\hex{00AE} : zone initialisée de façon conventionnelle ;
  \item \hex{00AE}--\hex{00B4} : initialisée à la mise sous tension ;
  \item \hex{00B4}--\hex{00BE} : initialisée à la remise à zéro du BASIC ;
  \item \hex{00BE}--\hex{0552} : zone générale de travail du système.
\end{itemize}

\subsection{Vecteurs, RST et hooks}
{\small
\begin{longtable}{@{}>{\ttfamily}l>{\ttfamily}l>{\raggedright\arraybackslash}p{10.5cm}@{}}
\toprule
\normalfont Adresse & \normalfont Instr. & Rôle \\
\midrule\endhead
\bottomrule\endfoot
000 & C9 & RST 00 (point d'arrêt de la carte moniteur XP140F) \\
003 & & nombre de caractères imprimés $\times$ taille \\
004 & 50 & longueur de ligne de la X-710 (80) \\
005 & & taille des caractères de la X-710 \\
008 & JP F52F & RST 08 : comparaison de symbole \\
00C-00D & & début de la ROM 1 \\
010 & JP F537 & RST 10 : test de caractère \\
013-014 & & fin de la ROM 1 \\
018 & JP C9C7 & RST 18 : émission d'une commande vers le T6834 \\
020 & JP EE45 & RST 20 : comparaison de HL et DE \\
024-025 & & début de la ROM 2 \\
026-027 & & fin de la ROM 2 \\
028 & JP E88F & RST 28 : émission de A vers le périphérique ouvert \\
02B & & BREAKFLAG \\
02C & C9 & INTERRUPT C (clignotement du curseur) \\
030 & JP FC2F & RST 30 : test du type de donnée \\
034 & JP C7A3 & INTERRUPT B (interruptions série) \\
038 & JP E906 & RST 38 : calcul d'adresse dans un tableau \\
03C & JP C799 & INTERRUPT A (interruptions du T6834) \\
03F & JP CED6 & appel de GPRPUT (imprimante parallèle) \\
042 & JP E512 & émission de A vers KBD, CASI \\
045-046 & & pointeur de la table du RST 38 \\
047-056 & & table des hooks \\
057 & JP E327 & minuterie automatique \\
05A & JP C16C & retour \\
05D & JP C1FA & \\
060 & JP C181 & \\
063 & JP C219 & \\
066 & C9 & NMI \\
069 & JP C582 & \\
06C & JP C0A1 & affichage ou masquage du curseur \\
06F & JP C210 & \\
072 & JP C231 & affiche C à la position du curseur \\
075 & JP C250 & SCREEN / image miroir \\
078 & JP C154 & \\
07B & JP C138 & \\
07E & JP C120 & \\
081 & JP C24B & TXTTAB \\
084 & JP F1AA & SN ERROR \\
087 & JP E403 & SCREEN \\
08A & JP F1AA & COLOR (X-720 Péritel), sinon SN ERROR \\
08D & JP CDF4 & PSET \\
090 & JP CDF7 & PRESET \\
093 & JP CE05 & POINT \\
096 & JP CE19 & LINE \\
099 & JP F1AA & PAINT (X-720 Péritel), sinon SN ERROR \\
09C & JP CE32 & CIRCLE \\
09F & JP C1BE & sortie CONSOLE (LCD par défaut) \\
0A2 & JP C90A & entrée CONSOLE (clavier par défaut) \\
0A5 & JP C16D & \\
0A8 & C9 & hook ABORT (JP 7535 avec la carte XP140F) \\
0AB & C9 & hook ERREUR (JP 756C avec la carte XP140F) \\
\end{longtable}
}
La zone \hex{057}--\hex{09E} relie l'afficheur LCD et l'afficheur vidéo
Péritel (X-720).

\nota{Pour le RST 18, la table indique \texttt{JP C9C7}, alors que le §\,9.5.1
donne \hex{C937}. Vérifié dans la ROM : au démarrage à froid, la routine
\hex{C3D8} copie \hex{AE} octets de \hex{C6DB} vers \hex{0000}
(\texttt{LDIR}) ; on y lit \texttt{C3 C7 C9} à la position \hex{18}.
\textbf{RST 18 saute donc en \hex{C9C7}}, et \hex{C937} tombe au milieu
d'une instruction. Le rôle indiqué est en revanche faux : \hex{C9C7} ne
s'adresse pas au T6834, elle donne le \textbf{signe de l'accumulateur
flottant} (\hex{044E}) : A = 0 et Z = 1 si la valeur est nulle, A = 1 si
elle est positive, A = \hex{FF} si elle est négative.}

\subsection{Variables du système}
{\small
\begin{longtable}{@{}>{\ttfamily}l>{\raggedright\arraybackslash}p{12.5cm}@{}}
\toprule
\normalfont Adresse & Contenu \\
\midrule\endhead
\bottomrule\endfoot
0B1 & indicateur RESUME en cours \\
0B2-0B3 & TXTTAB : début du programme BASIC \\
0B8 / 0B9 & position du curseur : X / Y \\
0BA & longueur d'une ligne de l'écran LCD \\
0BB / 0BC & première / dernière ligne de défilement \\
0BD & nombre de lignes de défilement \\
0D5-1D4 & tampon d'entrée \\
1D8 & indicateur de variable (0 : elle existe ; FF : à créer) \\
1D9 & type de donnée : 02 entier (\%), 03 chaîne (\$), 04 simple précision (!), 08 double précision (\#) \\
1DB-1DC & dernière ligne exécutée \\
1DD-1DE & STKTOP : fin de la pile \\
1DF-1E0 & MEMSIZ : fin de la zone BASIC \\
1E1-1E2 & prochaine entrée libre de la table des chaînes \\
1E3-200 & table de manipulation des chaînes \\
201 & longueur de la chaîne (0 à 255) \\
202-203 & adresse de la chaîne \\
204-205 & FRETOP : prochaine adresse libre de la zone des chaînes \\
206-207 & pointeur du dernier code exécuté / taille des tableaux / indicateur de PRINT USING \\
208-209 & D43A \\
20A-20B & pile des boucles FOR-NEXT \\
20E & indicateur de l'instruction FOR \\
20F & indicateur utilisé par C977 \\
210-211 & RAMSTR : début de la zone des fichiers en RAM \\
212-213 & RAMEND : fin de la zone des fichiers en RAM \\
214-263 & image miroir de l'écran LCD (80 octets = 20 $\times$ 4) \\
264-265 & sauvegarde de HL (pointeur BASIC) lors d'un SLEEP \\
266-267 & sauvegarde de SP lors d'un SLEEP \\
268-26B & mot-clé (jour, heure, minute, seconde) placé en RAMSTR$-4$ et RAMEND$-1$ \\
26C & INTIMAG : image du port F0 \\
26D & nombre d'octets de réponse attendus du T6834 \\
26E- & 1\textsuperscript{er} tampon de communication avec le T6834 \\
2C3-2C4 & adresse du pilote de sortie \\
2C5-2E8 & KBNTAB : table de INIT\# (5 $\times$ 8 octets : adresse + sauvegarde de HL, DE, BC) \\
2F5 & numéro du périphérique ouvert \\
2F6-2F7 & adresse du pilote \\
2F8-2FD & indicateurs de ALM\$, TIME\$ et DATE\$ (DATE\$ : 0 si actif, 1 si ALM\$ actif) \\
2FE & indicateur de CLOAD / CLOAD? \\
2FF-304 & nom du fichier cassette cherché (6 octets) \\
305-30A & nom du fichier lu sur la bande (SKIP ou FOUND) \\
30B & UCRIMAG \\
30C & RSERFLAG \\
30D & 1 si un INPUT est en cours, 0 sinon \\
30E & indicateur de l'instruction READ \\
30F-310 & appel à l'adresse 0B2 par RUN \\
311-312 & numéro de la ligne en cours d'exécution \\
313-314 & sauvegarde de SP pendant que l'interpréteur travaille \\
317-318 & numéro de la ligne où l'erreur s'est produite \\
319-31A & adresse du traitement d'erreur (ON ERROR) \\
31B & FF si un ON ERROR est actif, 0 sinon \\
31C-31D & adresse du point décimal dans le tampon de PRINT \\
31E-31F & adresse de reprise pour CONT \\
320-321 & adresse du dernier code exécuté lors d'une erreur \\
322-323 & VARTAB : table des variables \\
324-325 & ARRTAB : table des tableaux \\
326-327 & STREND : début de la zone libre \\
328-329 & pointeur des DATA après un RESTORE $n$ \\
32A-343 & types des variables par défaut (26 lettres) \\
344-345 & PRMSTK \\
348-3AB & pile (appel par B07C) \\
3AC-3AD & PRMPRV \\
414 & indicateur de recherche de variable \\
415-416 & adresse de la variable recherchée \\
41E- & tampon interne de PRINT \\
44E-455 & accumulateur flottant (44E : exposant ; 450-451 : entier) \\
456 & octet de travail de l'accumulateur \\
49F-4A6 & accumulateur 2 \\
4A7 & octet de travail de l'accumulateur 2 \\
4AF- & tampon de RND \\
4C6-4C9 & anciennes coordonnées X, Y de LINE \\
4CA-4CD & nouvelles coordonnées X, Y de LINE ( )--( ) \\
507- & 2\textsuperscript{e} tampon de réponse du T6834 \\
54D & indicateur de présence de la carte XP140F / fonction LOG \\
54E-54F & début de la zone de travail de la carte moniteur XP140F \\
\end{longtable}
}

\nota{Vérifié dans la ROM : \textbf{\hex{0B8} contient la ligne et \hex{0B9}
la colonne}, toutes deux comptées \textbf{à partir de 1}. \texttt{LOCATE x,y}
(\hex{E212}) charge H = x+1 et L = y+1 puis exécute
\texttt{LD (\$00B8),HL} (\hex{C164}) ; \texttt{CSRLIN} renvoie
(\hex{0B8})$-1$ et \texttt{POS} (\hex{0B9})$-1$ ; un retour chariot
remet \hex{0B9} à 1. Les programmes du livre qui écrivent
\texttt{LD HL,\$0502 : LD (\$00B8),HL} placent donc le curseur en colonne 5,
ligne 2.}

\nota{Vérifié dans la ROM : à chaque \texttt{RUN} (\hex{D234}), la ROM remplit
les 26 octets \hex{032A}--\hex{0343} avec \texttt{LD (HL),\$08} : \textbf{une
variable sans suffixe est en double précision} (14 chiffres), comme sur les
MSX, et non en simple précision. \texttt{PRINT 1/3} affiche donc
\texttt{.33333333333333}.}

\paragraph{BREAKFLAG (\hex{002B}).} Bit 0 : BREAK ; bit 1 : piles épuisées ;
bit 3 : pile de la carte épuisée ; bit 6 : arrêt par la minuterie ;
bit 7 : OFF.

\nota{Vérifié dans la ROM : à chaque instruction, l'interpréteur appelle
\hex{C000}, qui traite le bit 7 (extinction, par \hex{C5A6}), le
\textbf{bit 2 (pause} jusqu'à l'appui sur une touche, non cité par le
livre), les bits 1 et 3 (messages de piles faibles, par \hex{C06C}) puis le
bit 0 (BREAK : \hex{C02D} efface les bits 0 et 2 par \texttt{AND \$FA} et
vide le tampon du clavier). Le bit 6 n'est pas testé par ces routines.}

%======================================================================
\section{Principales routines de la ROM \pg{96-110}}\label{sec:routines}

Pour chaque routine, le livre donne l'entrée (registres à charger),
l'action, la sortie et les registres modifiés. Dans ce chapitre,
« accumulateur » désigne l'accumulateur flottant du BASIC (\hex{044E}) et
non le registre \reg{A} du Z80. Le livre emploie souvent « A » pour les
deux : le sens se déduit du contexte. Avant d'appeler une routine
arithmétique, il faut positionner l'indicateur de type \hex{01D9}
(2 entier, 4 simple précision, 8 double précision).

\subsection{Entrées/sorties}
\subsubsection*{Buzzer}
\begin{routines}
C2DF & IY = tonalité, H = durée, L = 0 & bip (comme BEEP) & -- & AF BC DE HL IY, EI \\
C331 & IY = tonalité, HL = durée, C = FF marche / 00 arrêt & commande le haut-parleur & -- & AF BC DE HL, EI \\
\end{routines}

\subsubsection*{Clavier}
\begin{routines}
C8C5 & -- & lit le clavier en affichant le curseur & A = caractère ; CY = 1 si interruption & AF \\
C90A & -- & lit un octet au clavier (comme INKEY\$) & A = caractère ; Z = 1 si aucune touche & AF \\
\end{routines}

\subsubsection*{Écran LCD}
\begin{routines}
C1BE & A = caractère & écrit à l'écran et dans l'image mémoire ; le curseur avance ; codes de contrôle admis & -- & aucun \\
C231 & H = X, L = Y, C = caractère & écrit C à la position indiquée ; le curseur ne bouge pas & -- & AF BC DE \\
CE9E & -- & efface l'écran (CLS) & -- & AF \\
C18A & A = code ASCII & affiche au curseur ; le curseur avance & -- & aucun \\
C0A1 & A = 0 ou 1 & 0 : affiche le curseur ; 1 : le cache & -- & AF BC DE \\
\end{routines}

\subsubsection*{Imprimante}
\begin{routines}
CFB7 & -- & initialise l'imprimante parallèle (mode texte, taille 2, couleur 0) & -- & AF BC DE HL \\
CFB0 & -- & envoie LF + CR à l'imprimante & -- & AF BC DE HL \\
CEF7 & A = octet & envoie A à l'imprimante (teste BUSY et BREAK) & -- & AF BC DE HL \\
CF11 & A = caractère & envoie A sur le port parallèle & -- & AF B E \\
\end{routines}

\subsubsection*{Sous-processeur T6834}
\begin{routines}
E428 & A = commande & envoie une commande sans paramètre & -- & AF C \\
C92F & A = commande, HL = paramètres, B = nombre de paramètres, C = octets de réponse, DE = tampon (26E) & envoie une commande avec paramètres & réponse en (DE) & AF BC HL ; bit 7 de A à 1 si B = 0 \\
C0BD & -- & vide les tampons du clavier et du sous-processeur & -- & AF \\
E348 & HL = adresse dans la RAM du T6834 & lit un octet & A = octet lu & AF BC DE \\
E334 & HL = adresse, A = donnée & écrit un octet dans la RAM du T6834 & -- & AF BC DE \\
\end{routines}

\subsubsection*{Routines générales}
\begin{routines}
C39D & -- & initialise les ports parallèle, série, optique, imprimante et cassette & -- & AF BC DE \\
CF2A & -- & attend que BUSY du port parallèle passe à l'état bas ; contrôle OFF, BREAK et les piles & -- & AF \\
\end{routines}

\subsection{Conversions}
\begin{routines}
CAE0 & accumulateur & flottant $\to$ entier (CINT) & -- & tous \\
CB08 & accumulateur & $\to$ simple précision (CSNG) & -- & tous \\
CB1E & entier dans l'accumulateur & entier $\to$ simple précision & -- & tous \\
CB90 & accumulateur & $\to$ double précision (CDBL) & -- & tous \\
F595 & HL = chaîne & ASCII $\to$ entier (s'arrête au premier caractère non numérique) & DE = résultat & \\
BA21 & HL = chaîne & ASCII $\to$ flottant ; indicateurs ajustés & accumulateur & \\
BB98 & HL = valeur & HL $\to$ ASCII, affiché à l'écran & écran & \\
BE5F & entier dans l'accumulateur, HL = tampon, B et C $>6$ & entier $\to$ ASCII & chaîne dans le tampon & tous \\
BBAB & accumulateur = flottant, B = chiffres avant la virgule, C = chiffres après, A = format & flottant $\to$ ASCII & tampon (HL), en général \hex{041E} & \\
\end{routines}
Format pour \hex{BBAB} : A\,=\,0 : conversion brute ; sinon bit 7 :
édition ; bit 6 : 1 ; bit 5 : « * » en tête ; bit 4 : « \$ » en tête ;
bit 3 : signe présent ; bit 2 : signe selon la valeur ; bit 0 : notation
exponentielle.

\nota{Vérifié dans la ROM : \hex{BBAB} commence par \texttt{XOR A} et ignore
donc A, B et C : c'est la conversion \textbf{en format libre} (FOUT), celle
qu'utilisent \texttt{PRINT} (\hex{EA2C}) et \texttt{STR\$} (\hex{D53D}). Elle
rend HL = début du texte, terminé par \hex{00}, avec une espace en tête pour
un nombre positif. Le format décrit ci-dessus correspond à une autre entrée
(\hex{BEEE}, \texttt{PRINT USING}).}

\subsection{Arithmétique}
\begin{routines}
CCC6 & HL, DE ; 1D9 = 2 & addition entière & HL (DE conservé) ; en cas de dépassement : accumulateur, type = 8 & \\
CCBB & HL, DE ; 1D9 = 2 & HL $-$ DE (entier) & idem & \\
CCE7 & HL, DE ; 1D9 = 2 & multiplication entière & idem & \\
FA79 & DE, HL & DE / HL (entier) & quotient en double précision dans l'accumulateur & \\
CDA9 & accumulateur, BC\string:DE ; 1D9 = 4 & addition en simple précision & accumulateur & \\
CDB2 & accumulateur = 1\textsuperscript{er}, BC\string:DE = 2\textsuperscript{e} & accumulateur $-$ BC\string:DE & accumulateur & \\
CDB7 & idem & accumulateur $\times$ BC\string:DE & accumulateur & \\
CDC2 & accumulateur = diviseur, BC\string:DE = dividende & BC\string:DE / accumulateur & accumulateur & \\
B20E & accumulateur, accumulateur 2 ; 1D9 = 8 & addition en double précision & accumulateur & \\
B200 & idem & accumulateur $-$ accumulateur 2 & accumulateur & \\
B376 & idem & multiplication & accumulateur & \\
B42F & idem & accumulateur / accumulateur 2 & accumulateur & \\
\end{routines}
\nota{Vérifié dans la ROM : les routines « simple précision » ne calculent pas
en simple précision. \hex{CDA9} range BC:DE dans l'accumulateur 2 (C en
\hex{049F}, B, E, D ensuite), complète l'accumulateur par des zéros
(\hex{CB98}) et saute dans \hex{B20E} : \textbf{le résultat est en double
précision} (type 8). L'arrondi à 6 chiffres n'a lieu qu'au rangement dans une
variable simple ou par \texttt{CSNG}. L'évaluateur choisit la routine dans une
table par type (\hex{F0FF} double, \hex{F10B} simple, \hex{F117} entier :
addition, soustraction, multiplication, division, puissance, comparaison),
l'opérande gauche ayant été empilé, puis saute en \hex{F9E4}. Pour la
comparaison en double (\hex{CAD9}), A vaut 1 si l'accumulateur est inférieur à
l'accumulateur 2, $-1$ s'il est supérieur, 0 s'ils sont égaux. Deux
comportements en découlent : un dépassement entier passe en double
(\texttt{PRINT 30000+30000} affiche 60000), et \texttt{CINT} (\hex{CAE0})
\textbf{tronque} au lieu d'arrondir (\texttt{CINT(2.6)} = 2,
\texttt{CINT(-2.6)} = $-2$).}

\subsubsection*{Comparaisons}
Dans ce tableau, A est le registre \reg{A} du Z80.
\begin{routines}
CA7B & BC\string:DE = 1\textsuperscript{er}, accumulateur = 2\textsuperscript{e} & comparaison en simple précision & A = $-1$ si BC\string:DE $>$ acc., 1 si $<$, 0 si $=$ & \\
CAA5 & HL, DE & comparaison entière & A = $-1$ si DE $>$ HL, 1 si $<$, 0 si $=$ & \\
CAD9 & accumulateur, accumulateur 2 & comparaison en double précision & A = $-1$ si acc.\,2 $>$ acc., 1 si $<$, 0 si $=$ & \\
C9FB & accumulateur ; 1D9 = 8 & signe en double précision & A = 1, $-1$ ou 0 & \\
\end{routines}

\subsection{Fonctions mathématiques}
\begin{routines}
C9DC & accumulateur & valeur absolue (ABS) & accumulateur & \\
CC23 & accumulateur & partie entière (INT) ; 1D9 actualisé & accumulateur & \\
\end{routines}

\subsection{Routines système}
\subsubsection*{Instructions RST}
\begin{center}\small
\begin{tabular}{@{}>{\ttfamily}l>{\ttfamily}l>{\raggedright\arraybackslash}p{11.5cm}@{}}
\toprule
\normalfont RST & \normalfont Routine & Rôle \\
\midrule
RST 00 & 0000 & \texttt{C9} (RET) : point d'arrêt de la carte moniteur XP140F \\
RST 08 & F52F & compare l'octet (HL) à l'octet qui suit le RST : s'ils sont égaux, HL est incrémenté et le programme reprend après cet octet ; sinon, « SN ERROR » \\
RST 10 & F537 & charge (HL) dans A et incrémente HL ; CY indique un caractère numérique ou alphabétique \\
RST 18 & C937 & envoie au T6834 la commande contenue dans A (sans paramètre) \\
RST 20 & EE45 & compare HL et DE : CY = 1 si HL $<$ DE ; CY = 0 si HL $>$ DE ; Z = 1 si égaux \\
RST 28 & E88F & envoie A au périphérique ouvert (en général le LCD) ; tous les registres sont conservés \\
RST 30 & FC2F & teste le type de la donnée (voir ci-dessous) \\
RST 38 & E906 & calcul de l'adresse d'un élément de tableau \\
\bottomrule
\end{tabular}
\end{center}
\nota{Vérifié dans la ROM : RST 18 saute en \hex{C9C7} et non en
\hex{C937}, et il ne commande pas le T6834 : il donne le signe de
l'accumulateur flottant (A = 0, 1 ou \hex{FF} ; Z = 1 si nul). Voir la note
de la table des vecteurs (§\,\ref{sec:zone}).}

\nota{Vérifié dans la ROM : RST 38 (\hex{E906}) n'est pas un calcul d'adresse
dans un tableau. Il est toujours suivi d'un \textbf{octet en ligne} n, que la
routine saute avant de revenir. Si n $<$ \hex{80}, il rend DE = mot lu en
HL$+$n (lecture dans une table). Si n $\ge$ \hex{80}, il appelle la
fonction n$-$\hex{80} du \textbf{pilote du périphérique courant} (table en
\hex{0045}) : \texttt{PRINT} (\texttt{FF 83}), la position et la largeur pour
\texttt{PRINT} et \texttt{TAB} (\hex{E8EE} : \texttt{FF 81}), le placement
pour \texttt{TAB} (\texttt{FF 8F}), \texttt{INIT\#}\dots{} Un désassemblage
qui ignore cet octet devient faux à partir de là.}

\begin{center}\small
\begin{tabular}{@{}clll@{}}
\toprule
\hex{1D9} & Type & Indicateurs & A \\
\midrule
2 & entier & NC, Z, M & $-1$ \\
3 & chaîne & Z, C, P & 0 \\
4 & simple précision & NZ, C, P & 1 \\
8 & double précision & NZ, NC, P & 5 \\
\bottomrule
\end{tabular}
\end{center}
\nota{Vérifié dans la ROM (exécution de \hex{FC2F}) : la ligne de l'entier est
fausse. Pour un entier, RST 30 rend \textbf{C = 1, Z = 0, M} (A = $-1$).
Ainsi, C = 1 pour l'entier, la chaîne et la simple précision, et C = 0 pour la
double précision : le répartiteur des fonctions (\hex{FBEB}) convertit en
double l'argument de \texttt{SQR}, \texttt{RND}, \texttt{LOG}, \texttt{EXP},
\texttt{COS}, \texttt{SIN}, \texttt{TAN} et \texttt{ATN} par
\texttt{CALL C,\$CB90}. Ces routines donnent un résultat faux si on leur passe
directement un entier.}

\subsubsection*{Transferts de données}
\begin{routines}
CA0B & accumulateur (simple précision) & empile l'accumulateur (1D9 non testé) & pile & \\
CA1B & BC, DE & BC\string:DE $\to$ accumulateur & accumulateur & \\
CA26 & accumulateur & accumulateur $\to$ BC\string:DE & BC, DE & \\
CA42 & accumulateur, HL = tampon & accumulateur $\to$ (HL) & tampon rempli & \\
CA51 & B = nombre d'octets, DE = source, HL = destination & copie un bloc & HL et DE croissants & \\
CA58 & idem & copie un bloc & HL et DE décroissants & \\
CAEF & HL & HL $\to$ accumulateur ; 1D9 = 2 & accumulateur & \\
CA5F & accumulateur 2 & accumulateur 2 $\to$ accumulateur & accumulateur & \\
CA67 & accumulateur & accumulateur $\to$ accumulateur 2 & accumulateur 2 & \\
\end{routines}

\subsection{Routines secondaires}
\begin{routines}
C02D & -- & lit BREAKFLAG ; si le bit 6 est à 1 : le remet à 0, vide le tampon des touches et met Z à 1 ; sinon retour avec Z positionné & -- & aucun \\
C5A6 & -- & contrôle des piles : bit 1 $\to$ « LOW BATTERY », bit 3 $\to$ « CARD LOW BATTERY » (Z = 0) ; sinon Z = 1 & A = BREAKFLAG & aucun \\
EE13 & HL = table de sauts, C = décalage & saut dans une table de vecteurs & HL modifié & AF BC DE HL \\
C557 & -- & lit le commutateur de la prise d'extension & HL = 4000 (RAM) ou 2000 (ROM) & AF HL \\
C5C3 & -- & coupe l'alimentation & -- & tous \\
C5CD & A = 1 ou 2 & 1 : valide le programme de lancement et éteint ; 2 : l'invalide et coupe l'alimentation & -- & aucun \\
E327 & -- & réarme la minuterie d'extinction (15 min) & -- & AF B \\
C62E & HL = adresse & teste si l'adresse est en RAM & Z = 1 si oui & AF B \\
C59B & -- & SLEEP : sauvegarde l'état des piles et HL (264), puis passe en SLEEP & -- & aucun \\
E8E8 & -- & lit BREAKFLAG & Z = 1 si un bit est à 1 & AF \\
D5B0 & HL = chaîne (terminée par 00 ou 0D) & affiche un message (utilise la zone des chaînes) & -- & \\
FEF7 & HL = chaîne (terminée par 00) & affiche un message, sans la zone des chaînes & -- & \\
EBF2 & -- & saisie au clavier avec « ? » & HL = tampon d'entrée $-1$ & \\
EBFF & -- & saisie au clavier sans « ? » & HL = tampon d'entrée $-1$ & \\
F30D & DE = numéro de ligne & cherche une ligne BASIC & si C = Z = 1 : BC = adresse de la ligne & \\
FE5E & A = expression & évaluation entière ($<$ 256) & DE, poids faible dans A ; « FC ERROR » si $>255$ & \\
FFCC & expression & évaluation entière & DE et A & \\
\end{routines}
\nota{Le livre décrit une seconde fois \hex{C5A6} (p.~108) avec un autre rôle
(remise à zéro de BREAKFLAG et extinction) : l'une des deux adresses est
probablement fausse. De même, pour \hex{C59B}, l'adresse de sauvegarde des
piles est imprimée « 226h ».}

\nota{Vérifié dans la ROM :
\begin{itemize}
  \item \hex{C5A6} correspond à la \emph{seconde} description du livre : si le
        bit 7 de BREAKFLAG (OFF) est à 1, il remet BREAKFLAG à 0 et coupe
        l'alimentation ; sinon il revient avec A = BREAKFLAG. Les messages
        « LOW BATTERY » (bits 1 et 3) sont affichés par \hex{C06C} ;
  \item \hex{C02D} teste le \textbf{bit 0} (BREAK) et non le bit 6 : s'il est
        à 1, il efface les bits 0 et 2 (\texttt{AND \$FA}) et vide les tampons
        du clavier et du T6834 (\hex{C0BD}) ;
  \item \hex{E8E8} fait \texttt{LD A,(\$002B) : AND \$83} : \textbf{Z = 0}
        (et non Z = 1) si BREAK, piles faibles ou OFF est signalé ;
  \item \hex{EBF2} et \hex{EBFF} rendent HL = \hex{00D4} ; la ligne saisie
        se termine par \hex{00} et \textbf{CY = 1 signale un appui sur BREAK}
        pendant la saisie (\texttt{INPUT} saute alors en \hex{D2CA}).
\end{itemize}}

\subsubsection*{Points d'entrée divers}
\begin{center}\small
\begin{tabular}{@{}>{\ttfamily}ll@{}}
\toprule
\normalfont Adresse & Rôle \\
\midrule
F680 & CONT général du système \\
C528 & affichage du message de copyright \\
C3C3 & initialisation générale (démarrage à froid) \\
FB4F & entrée d'une routine en hexadécimal \\
C4D1 & « CREATE SYSTEM ? » (Y, RETURN) : efface complètement une carte mémoire \\
F23D & mode d'attente avec curseur \\
F1AA & affiche « SN ERROR » \\
F1C7 & E = numéro d'erreur : affiche l'erreur correspondante \\
\bottomrule
\end{tabular}
\end{center}

%======================================================================
\section{Périphériques \pg{111-115}}

\subsection{Cartes mémoire}
Il existe deux types de cartes : RAM seule (4 ou 8\,Ko), ou ROM (8\,Ko en
général) avec 4\,Ko de RAM, dont 3\,Ko pour l'utilisateur (1\,Ko sert de
zone de travail à la carte).

\begin{center}\small
\begin{tabular}{@{}cl|cl@{}}
\toprule
Broche & Signal & Broche & Signal \\
\midrule
1 & VCC1 (alimentation du X-07) & 11--18 & D7 \dots{} D0 \\
2 & VCC2 (alimentation du X-07) & 19--31 & A12 \dots{} A0 \\
3 & NC & 32 & masse \\
4 & V+ (détection LOW BAT.) & 33 & masse \\
5 & NC & & \\
6 & accès RAM (RAM ON) & & \\
7 & sélection RAM (active à l'état bas) & & \\
8 & $\overline{\mbox{WR}}$ & & \\
9 & $\overline{\mbox{RD}}$ & & \\
10 & sélection ROM (active à l'état bas) & & \\
\bottomrule
\end{tabular}
\end{center}

\subsection{Imprimante graphique X-710}
La X-710 possède son propre microprocesseur 4 bits, une ROM de
4096 $\times$ 10 bits et une RAM de 256 $\times$ 4 bits, avec une horloge
d'environ 600\,kHz. L'octet de déplacement du stylo se décompose ainsi :
\begin{center}\small
\begin{tabular}{@{}ll@{}}
\toprule
Bits & Signification \\
\midrule
7 & 1 : fin de déplacement ; 0 : déplacement \\
4--6 & longueur du mouvement \\
3 & 0 : stylo levé ; 1 : stylo baissé (tracé) \\
0--2 & direction du mouvement \\
\bottomrule
\end{tabular}
\end{center}

\paragraph{Brochage du contrôleur (fig.~29, boîtier FPP 54 broches).}
{\small
\begin{longtable}{@{}cl>{\raggedright\arraybackslash}p{11cm}@{}}
\toprule
Broche & Nom & Fonction \\
\midrule\endhead
\bottomrule\endfoot
1 & D4 & SCS (échelle) : L $\to$ échelle 1 après reset, H $\to$ échelle 0 \\
2 & D5 & CRM : L $\to$ CR seul, H $\to$ CR + LF \\
3 & D6 & EXEN (validation de la ROM externe) \\
4--7 & D7--D10 & touches : avance papier, recul papier, changement de couleur, changement de stylo \\
8 & D11 & PWS (largeur du papier) : L 58\,mm, H 114\,mm \\
9 & D12 & CNS (nombre de couleurs) : L 1, H 4 \\
10 & D13 & HLTEN (Halt Enable) \\
11 & D14 & STR (reset de la bascule Strobe) \\
12 & D15 & BSR (reset de la bascule Busy) \\
13--16 & R40--R43 & données parallèles D0--D3 (code caractère) \\
17--20 & R50--R53 & données parallèles D4--D7 \\
21 & RESET & reset \\
22 & GND & alimentation négative \\
23--24 & OSC & oscillateur \\
25 & HLT & signal Halt \\
26 & TEST & relié à VCC \\
27 & VCC & alimentation positive \\
28--35 & R00--R13 & EXD0--EXD7 : données de la ROM externe (reliées à VCC si elle est absente) \\
36 & INT0 & entrée Strobe (front initial) \\
37 & INT1 & reliée à GND \\
38--41 & R20--R23 & RA0--RA3 : adresse de la ROM externe \\
42--45 & R30--R33 & $\overline{\mbox{XD}}$\dots$\overline{\mbox{XA}}$ : moteur X \\
46--49 & R60--R63 & $\overline{\mbox{YD}}$\dots$\overline{\mbox{YA}}$ : moteur Y \\
50 & D0 & $\overline{\mbox{PEN UP}}$ \\
51 & D1 & $\overline{\mbox{PEN DOWN}}$ \\
52 & D2 & $\overline{\mbox{COLOR DETECT}}$ \\
53 & D3 & $\overline{\mbox{BAS}}$ (strobe d'adresse de la ROM externe) \\
54 & NC & -- \\
\end{longtable}
}
Le schéma complet de la X-710 (fig.~30) se trouve p.~115 du livre :
contrôleur IC3, pilote de moteurs IC4, mécanisme DPG-23, alimentation
AD-5 et batterie NiCd de 6\,V.

\subsection{Interface vidéo X-720}
\begin{itemize}
  \item 6\,Ko de mémoire vidéo de \hex{8000} à \hex{97FF}, entièrement
        utilisés par les modes \texttt{SCREEN 5} et \texttt{6}. En basse
        résolution, deux pages d'écran sont disponibles, ce qui permet de
        basculer très vite de l'une à l'autre.
  \item 4\,Ko de ROM.
  \item La X-720 redirige les hooks de la zone système (comme la carte
        XP140F) pour ajouter SCREEN, PAINT et COLOR, et pour étendre LINE,
        PSET, PRESET et CIRCLE.
  \item Le connecteur à l'arrière était prévu pour des cartouches ROM ou des
        extensions de RAM, mais aucune n'a été commercialisée.
\end{itemize}
\nota{Le livre écrit « 800h et 97FFh » pour la mémoire vidéo (il faut lire
\hex{8000}) et « 4000h à B000h » pour la ROM, alors que la carte mémoire
(fig.~19) place la ROM TV en \hex{A000}--\hex{AFFF}.}

\newpage
%======================================================================
\part*{3\textsuperscript{e} partie du livre : Applications \pg{117-135}}
\addcontentsline{toc}{part}{\texorpdfstring{3\textsuperscript{e}}{3e} partie du livre : Applications}

La 3\textsuperscript{e} partie du livre met en pratique les notions des deux
premières : huit petites routines en langage machine Z80, chacune avec son
principe, son listing assembleur, un organigramme et un programme BASIC qui
charge les octets en mémoire par \texttt{POKE}. Les listings ont été
vérifiés en désassemblant les octets des lignes \texttt{DATA} ; les erreurs
du livre y sont corrigées.

Les routines de la ROM utilisées sont décrites à la section~\ref{sec:routines}.

%======================================================================
\section{Vue d'ensemble}

\subsection{Les huit programmes}
\begin{center}\small
\begin{tabular}{@{}cl>{\ttfamily}l>{\raggedright\arraybackslash}p{6.6cm}@{}}
\toprule
Fig. & Programme & \normalfont Implantation & Ce qu'il fait \\
\midrule
31 & Affichage & 1B00--1B24 & dessine un motif de 8$\times$8 points en écrivant dans la RAM LCD du T6834 \\
32 & Inversion vidéo & 1C00--1C15 & inverse la mémoire vidéo de la X-720 (SCREEN 5 et 6) \\
33 & Inversion LCD & 1C00--1C1C & inverse l'écran LCD \\
34 & Affichage d'un titre & 1C00--1C2C & fige un titre sur la première ligne de l'écran \\
35 & Copyright personnel & 2000--2027 & affiche un message à la mise sous tension (mot-clé \texttt{love}) \\
36 & Écriture sur la X-710 & 1C00--1C26 & imprime un message sur la table traçante X-710 \\
37 & Bruitages & 1C00--1C2B & produit des sons par le générateur de bauds \\
38 & Redéfinition de touche & 1C00--1C20 & associe un message à une touche \\
\bottomrule
\end{tabular}
\end{center}

Les programmes sont indépendants. Sauf le copyright, ils s'implantent tous en
\hex{1B00} ou en \hex{1C00}, à la fin des 8\,Ko de RAM interne : il faut
donc les charger un par un. La zone choisie doit rester hors du programme
BASIC chargeur. Si le programme BASIC est long, on peut réserver la mémoire
avec \texttt{CLEAR}.

\subsection{Le chargeur BASIC type}
Chaque programme est fourni avec un chargeur BASIC sur le même modèle :
\begin{verbatim}
10 CLS:PRINT"un instant !"
20 FOR I=&H1C00 TO &H1C15            ' première et dernière adresse
30 READ A$:POKE I,VAL("&H"+A$)       ' un octet hexadécimal par DATA
40 NEXT I
50 EXEC &H1C00                       ' lance la routine
60 DATA ...
\end{verbatim}
La routine se termine par \texttt{RET} et rend la main au BASIC. Dans les
DATA, les octets sont donnés en hexadécimal, sans préfixe.

\subsection{Routines de la ROM utilisées}
\begin{center}\small
\begin{tabular}{@{}>{\ttfamily}l>{\raggedright}p{8.5cm}l@{}}
\toprule
\normalfont Adresse & Rôle & Programmes \\
\midrule
CE9E & efface l'écran (CLS) & titre, copyright \\
E334 & écrit A dans la RAM du T6834, à l'adresse HL & affichage, LCD, titre \\
E348 & lit dans A l'octet de la RAM du T6834 à l'adresse HL & LCD \\
FEF7 & affiche la chaîne pointée par HL (terminée par \hex{00}) & titre, copyright \\
D5B0 & affiche la chaîne pointée par HL (zone des chaînes BASIC) & redéfinition \\
F23D & retour au mode « attente curseur » & copyright, redéfinition \\
C1BE & sortie d'un caractère sur la console (LCD) & redéfinition \\
CFB7 & initialise l'imprimante & X-710 \\
CED6 & envoie A à l'imprimante (GPRPUT) & X-710 \\
\bottomrule
\end{tabular}
\end{center}

%======================================================================
\section{Affichage d'un motif \texorpdfstring{8$\times$8}{8x8} \pg{119-120}}

\paragraph{But.} Dessiner un motif de 8$\times$8 points en écrivant
directement dans la mémoire de l'écran LCD, qui appartient au T6834. On n'est
plus limité aux caractères de 6$\times$8 points. Le principe est le même que
celui de l'instruction BASIC \texttt{FONT\$}.

\paragraph{Principe.} Dans la RAM de l'écran, la plus petite unité
adressable est une « ligne » de 8 points horizontaux : un bit à 1 est un
point allumé, un bit à 0 un point éteint. Le motif est donc décrit par une
suite de 8 octets, un par ligne. La routine écrit ces octets l'un après
l'autre dans la RAM du T6834 avec \hex{E334}, en commençant à \hex{E000}
(coin haut gauche de l'écran). Elle ajoute \hex{10} à l'adresse après chaque
ligne, ce qui correspond à une ligne de points plus bas.

\fiche{\hex{1B00}--\hex{1B1A} (code) ; \hex{1B1B}--\hex{1B24} (données)}%
      {27 octets de code + 10 octets de données}{\texttt{EXEC \&H1B00}}

\begin{asm}
1B00 & 06 08    &     & LD B,\$08      & nombre de lignes \\
1B02 & 21 1B 1B &     & LD HL,TB       & adresse de la table du motif \\
1B05 & 7E       & LO  & LD A,(HL)      & octet du motif \\
1B06 & E5       &     & PUSH HL        & \\
1B07 & 2A 23 1B &     & LD HL,(AD)     & adresse d'écriture dans la RAM LCD \\
1B0A & C5       &     & PUSH BC        & \\
1B0B & CD 34 E3 &     & CALL \$E334    & écriture dans la RAM du T6834 \\
1B0E & C1       &     & POP BC         & \\
1B0F & 11 10 00 &     & LD DE,\$0010   & une ligne de points plus bas \\
1B12 & 19       &     & ADD HL,DE      & \\
1B13 & 22 23 1B &     & LD (AD),HL     & \\
1B16 & E1       &     & POP HL         & \\
1B17 & 23       &     & INC HL         & octet suivant du motif \\
1B18 & 10 EB    &     & DJNZ LO        & \\
1B1A & C9       &     & RET            & \\
1B1B & FF 81 \dots & TB & DEFB \$FF,\$81,\$81,\$81, & motif : un carré \\
     & \dots\ 81 FF &   & \phantom{DEFB }\$81,\$81,\$81,\$FF & \\
1B23 & 00 E0    & AD  & DEFB \$00,\$E0 & adresse de départ : \hex{E000} \\
\end{asm}

\paragraph{Le motif.} Chaque octet code une ligne du dessin :
\begin{center}\small\ttfamily
\begin{tabular}{@{}lll@{}}
\toprule
\normalfont Dessin & \normalfont Binaire & \normalfont Hexa \\
\midrule
\#\#\#\#\#\#\#\# & 11111111 & FF \\
\#......\# & 10000001 & 81 \\
\#......\# & 10000001 & 81 \\
\#......\# & 10000001 & 81 \\
\#......\# & 10000001 & 81 \\
\#......\# & 10000001 & 81 \\
\#......\# & 10000001 & 81 \\
\#\#\#\#\#\#\#\# & 11111111 & FF \\
\bottomrule
\end{tabular}
\end{center}
Pour un autre dessin, il suffit de remplacer les 8 octets de \hex{1B1B} à
\hex{1B22}. On change l'emplacement en modifiant les deux octets de
\hex{1B23} (adresse dans la RAM LCD, poids faible en premier).

\paragraph{Chargeur BASIC.}
\begin{verbatim}
10 CLS:PRINT"un instant !"
20 FOR I=&H1B00 TO &H1B24
30 READ A$:POKE I,VAL("&H"+A$)
40 NEXT I
50 DATA 06,08,21,1B,1B,7E,E5,2A,23,1B,C5,CD,34,E3,C1,11,10,00,19,22,23,1B
60 DATA E1,23,10,EB,C9,FF,81,81,81,81,81,81,FF,00,E0
\end{verbatim}
Le chargeur ne lance pas la routine : taper \texttt{EXEC \&H1B00}.

\nota{La routine modifie l'adresse stockée en \hex{1B23} : pour la relancer
au même endroit, il faut d'abord remettre \hex{00},\hex{E0} dans ces deux
octets (ou relancer le chargeur). Le texte renvoie au « paragraphe 7.7 »,
qui n'existe pas : il s'agit du §\,7.6.2 (RAM du T6834).}

%======================================================================
\section{Inversion de l'écran vidéo \pg{121-122}}

\paragraph{But.} Inverser les couleurs de l'écran de l'interface vidéo X-720
de façon quasi instantanée, dans les modes \texttt{SCREEN 5} et \texttt{6}.

\paragraph{Principe.} On utilise un \texttt{XOR} avec \hex{FF} : $1 \oplus 1
= 0$ et $0 \oplus 1 = 1$, donc chaque bit est inversé. La routine parcourt
toute la mémoire vidéo, de \hex{8000} à \hex{97FF}, octet par octet. Cette
mémoire est directement accessible au Z80, d'où la rapidité.

\fiche{\hex{1C00}--\hex{1C15}}{22 octets}{\texttt{EXEC \&H1C00}}

\begin{asm}
1C00 & 06 FF    &    & LD B,\$FF      & masque d'inversion \\
1C02 & 21 00 80 &    & LD HL,\$8000   & début de la mémoire vidéo \\
1C05 & 7E       & 1E & LD A,(HL)      & \\
1C06 & A8       &    & XOR B          & inversion \\
1C07 & 77       &    & LD (HL),A      & \\
1C08 & 7C       &    & LD A,H         & \\
1C09 & FE 97    &    & CP \$97        & poids fort de l'adresse de fin ? \\
1C0B & 28 03    &    & JR Z,FI        & \\
1C0D & 23       & 2E & INC HL         & \\
1C0E & 18 F5    &    & JR 1E          & \\
1C10 & 7D       & FI & LD A,L         & \\
1C11 & FE FF    &    & CP \$FF        & poids faible de l'adresse de fin ? \\
1C13 & C8       &    & RET Z          & fin à \hex{97FF} \\
1C14 & 18 F7    &    & JR 2E          & \\
\end{asm}

\paragraph{Chargeur BASIC et démonstration.} Le programme passe en
\texttt{SCREEN 5}, trace un trait, puis inverse l'écran en boucle.
\begin{verbatim}
 10 CLS:PRINT"un instant !"
 20 FOR I=&H1C00 TO &H1C15
 30 READ A$:POKE I,VAL("&H"+A$)
 40 NEXT
 50 DATA 06,FF,21,00,80,7E,A8,77,7C,FE,97,28,03,23,18,F5,7D,FE,FF,C8,18,F7
 60 SCREEN 5
 70 PRINT"INVERSE VIDEO"
 80 FOR I=0 TO 100
 90 PSET (I,50):NEXT
100 EXEC&H1C00
110 FOR I=0 TO 500:NEXT:GOTO 100
\end{verbatim}
Ce programme nécessite l'interface X-720.

%======================================================================
\section{Inversion de l'écran LCD \pg{123-124}}

\paragraph{But.} Même chose pour l'écran LCD : inverser chaque octet de sa
mémoire par un \texttt{XOR} avec \hex{FF}.

\paragraph{Principe.} La mémoire de l'écran LCD n'est pas adressée par le
NSC800 mais par le T6834. On passe donc par les routines \hex{E348}
(lecture) et \hex{E334} (écriture) de la ROM, octet par octet, de
\hex{E000} à \hex{E1FE}. Les octets \hex{E00F}, \hex{E01F} et \hex{E02F}
n'ont pas besoin d'être inversés. Les auteurs les inversent quand même :
les tester prendrait plus de temps que de les traiter.

\fiche{\hex{1C00}--\hex{1C1C}}{29 octets}{\texttt{EXEC \&H1C00}}

\begin{asm}
1C00 & 21 FF DF &    & LD HL,\$DFFF   & début de la RAM LCD $-1$ \\
1C03 & 23       & 1E & INC HL         & \\
1C04 & CD 48 E3 &    & CALL \$E348    & lecture dans la RAM LCD \\
1C07 & 06 FF    &    & LD B,\$FF      & \\
1C09 & A8       &    & XOR B          & inversion \\
1C0A & CD 34 E3 &    & CALL \$E334    & écriture dans la RAM LCD \\
1C0D & E5       &    & PUSH HL        & \\
1C0E & 7C       &    & LD A,H         & \\
1C0F & FE E1    &    & CP \$E1        & poids fort de l'adresse de fin ? \\
1C11 & 28 03    &    & JR Z,FI        & \\
1C13 & E1       &    & POP HL         & \\
1C14 & 18 ED    &    & JR 1E          & \\
1C16 & 7D       & FI & LD A,L         & \\
1C17 & FE FE    &    & CP \$FE        & poids faible de l'adresse de fin ? \\
1C19 & E1       &    & POP HL         & \\
1C1A & C8       &    & RET Z          & fin à \hex{E1FE} \\
1C1B & 18 E6    &    & JR 1E          & \\
\end{asm}

\paragraph{Chargeur BASIC.}
\begin{verbatim}
10 CLS:PRINT"un instant"
20 FOR I=&H1C00 TO &H1C1C
30 READ A$:POKE I,VAL("&H"+A$)
40 NEXT
45 EXEC&H1C00
50 DATA 21,FF,DF,23,CD,48,E3,06,FF,A8,CD,34,E3,E5,7C,FE,E1,28,03,E1,18,ED,7D
60 DATA FE,FE,E1,C8,18,E6
\end{verbatim}
\nota{Le livre écrit \texttt{FOR I=\&H1C00 TO \&H1C1B} : le dernier octet
(\hex{E6}, à \hex{1C1C}) n'est alors jamais écrit, et l'instruction
\texttt{JR} finale saute n'importe où. La ligne 20 est corrigée ici.}

\paragraph{Attention.} Sur certaines versions du X-07, des octets échappent
à l'inversion quand le T6834 est débordé par le flux de données. Dans ce
cas, on peut utiliser à la place la commande \hex{13} du T6834 (PEOR, qui
inverse un point).

%======================================================================
\section{Affichage d'un titre \pg{125-126}}

\paragraph{But.} Figer un message sur la première ligne de l'écran. Ce
« titre » ne disparaît pas quand l'écran défile ; seul un \texttt{CLS}
l'efface.

\paragraph{Principe.} Après un effacement de l'écran (\hex{CE9E}), la
routine écrit dans la RAM du T6834 les limites de la zone de défilement :
ligne de début 1 en \hex{C09E}, ligne de fin 3 en \hex{C096}. La ligne 0
n'est donc plus concernée par le défilement. Le message est ensuite affiché
par \hex{FEF7}. Il doit se terminer par \hex{00}.

\fiche{\hex{1C00}--\hex{1C19} (code) ; \hex{1C1A}--\hex{1C2C} (message)}%
      {26 octets de code + 19 octets de message}{lancé par le chargeur}

\begin{asm}
1C00 & CD 9E CE &    & CALL \$CE9E    & efface l'écran \\
1C03 & 3E 01    &    & LD A,\$01      & début du défilement : ligne 1 \\
1C05 & 21 9E C0 &    & LD HL,\$C09E   & adresse dans la RAM du T6834 \\
1C08 & CD 34 E3 &    & CALL \$E334    & \\
1C0B & 3E 03    &    & LD A,\$03      & fin du défilement : ligne 3 \\
1C0D & 21 96 C0 &    & LD HL,\$C096   & \\
1C10 & CD 34 E3 &    & CALL \$E334    & \\
1C13 & 21 1A 1C &    & LD HL,ME       & \\
1C16 & CD F7 FE &    & CALL \$FEF7    & affichage du titre \\
1C19 & C9       &    & RET            & \\
1C1A & 20 2A \dots & ME & DEFM " *** PROGRAMME ***" & message \\
1C2C & 00       &    & DEFB \$00      & fin du message \\
\end{asm}

\paragraph{Chargeur BASIC.}
\begin{verbatim}
10 CLS:PRINT"un instant !"
20 FOR I=&H1C00 TO &H1C2C
30 READ A$:POKE I,VAL("&H"+A$)
40 NEXT I:EXEC&H1C00
50 DATA CD,9E,CE,3E,01,21,9E,C0,CD,34,E3,3E,03,21,96,C0,CD,34,E3,21,1A,1C
60 DATA CD,F7,FE,C9,20,2A,2A,2A,20,50,52,4F,47,52,41,4D,4D,45,20,2A,2A,2A,00
\end{verbatim}
Pour un autre titre, remplacer les codes ASCII de la ligne 60 (20 caractères
au plus, terminés par \hex{00}) et ajuster la borne de la ligne 20.

\nota{Le livre annonce 44 octets, de \hex{1C00} à \hex{1C2B}. Avec le
\hex{00} final, le programme occupe en réalité 45 octets, jusqu'à
\hex{1C2C}, comme l'indique la ligne 20.}

%======================================================================
\section{Copyright personnel \pg{127-128}}

\paragraph{But.} Afficher un message personnel à chaque mise sous tension
du X-07.

\paragraph{Principe.} À l'initialisation, la ROM cherche le mot-clé
\texttt{love} tous les \hex{800} octets à partir de \hex{2000} (voir
§\,7.4 du livre). Les 6 octets qui suivent le mot-clé sont trois adresses :
les deux premières sont appelées au démarrage, la troisième après un
\texttt{SLEEP}. Ici, la première adresse pointe vers la routine
d'affichage. Les deux autres pointent vers un simple saut au mode
« attente curseur ».

\paragraph{Emplacement.} La routine doit commencer à une adresse testée par
la ROM : \hex{2000}, \hex{2800}, \hex{3000}, \hex{3800}, \hex{4000},
\hex{4800} ou \hex{5000}. Il faut donc de la RAM à cet endroit (carte RAM
ou extension de 8\,Ko). La routine est relogeable, mais les trois adresses
qui suivent \texttt{love} doivent être modifiées en conséquence. Par
exemple, à partir de \hex{3000} :
\begin{verbatim}
        ORG  $3000
        DEFM love
        DEFB $0A,$30
        DEFB $19,$30
        DEFB $19,$30
\end{verbatim}

\fiche{\hex{2000}--\hex{2027} (en-tête \hex{2000}--\hex{2009}, code
\hex{200A}--\hex{201B}, message \hex{201C}--\hex{2027})}%
      {40 octets}{mettre le X-07 hors tension puis sous tension}

\begin{asm}
2000 & 6C 6F 76 65 &    & DEFM love      & mot-clé cherché par la ROM \\
2004 & 0A 20    &    & DEFB \$0A,\$20 & 1\textsuperscript{er} sous-programme : \hex{200A} \\
2006 & 19 20    &    & DEFB \$19,\$20 & 2\textsuperscript{e} sous-programme : \hex{2019} \\
2008 & 19 20    &    & DEFB \$19,\$20 & sous-programme après SLEEP : \hex{2019} \\
200A & CD 9E CE &    & CALL \$CE9E    & efface l'écran \\
200D & 21 02 03 &    & LD HL,\$0302   & X = 2, Y = 3 \\
2010 & 22 B8 00 &    & LD (\$00B8),HL & position du curseur \\
2013 & 21 1C 20 &    & LD HL,ME       & \\
2016 & CD F7 FE &    & CALL \$FEF7    & affichage du message \\
2019 & C3 3D F2 &    & JP \$F23D      & attente curseur \\
201C & 76 69 \dots & ME & DEFM "vive le X07" & message \\
2027 & 00       &    & DEFB \$00      & fin du message \\
\end{asm}

\paragraph{Chargeur BASIC.} Le programme éteint le X-07 (\texttt{OFF}) :
le message s'affiche au rallumage.
\begin{verbatim}
10 CLS:PRINT"un instant"
20 FOR I=&H2000 TO &H2027
30 READ A$:POKE I,VAL("&H"+A$)
40 NEXT I
50 OFF
60 DATA 6C,6F,76,65,0A,20,19,20,19,20,CD,9E,CE,21,02,03,22,B8,00,21,1C,20
70 DATA CD,F7,FE,C3,3D,F2,76,69,76,65,20,6C,65,20,58,30,37,00
\end{verbatim}

\nota{Le livre annonce une routine de 29 octets, de \hex{200A} à
\hex{2026}. Elle va en réalité jusqu'à \hex{2027} (30 octets, une
correction manuscrite du lecteur sur le scan l'indique aussi), soit 40
octets avec l'en-tête.}

%======================================================================
\section{Écriture sur la table traçante X-710 \pg{129-130}}

\paragraph{But.} Imprimer un message sur la X-710. L'intérêt est surtout de
montrer comment utiliser les routines de l'imprimante.

\paragraph{Principe.} L'imprimante est initialisée par \hex{CFB7}. Les
caractères du message sont ensuite envoyés un par un à \hex{CED6}, jusqu'au
code \hex{00} qui marque la fin. Les registres sont sauvegardés autour de
chaque appel.

\fiche{\hex{1C00}--\hex{1C17} (code) ; \hex{1C18}--\hex{1C26} (message)}%
      {24 octets de code + 15 octets de message}{\texttt{EXEC \&H1C00}}

\begin{asm}
1C00 & CD B7 CF &    & CALL \$CFB7    & initialisation de l'imprimante \\
1C03 & 21 18 1C &    & LD HL,TB       & \\
1C06 & 7E       & 1E & LD A,(HL)      & \\
1C07 & FE 00    &    & CP \$00        & fin de la chaîne ? \\
1C09 & C8       &    & RET Z          & \\
1C0A & F5       &    & PUSH AF        & \\
1C0B & C5       &    & PUSH BC        & \\
1C0C & D5       &    & PUSH DE        & \\
1C0D & E5       &    & PUSH HL        & \\
1C0E & CD D6 CE &    & CALL \$CED6    & envoi à l'imprimante \\
1C11 & E1       &    & POP HL         & \\
1C12 & D1       &    & POP DE         & \\
1C13 & C1       &    & POP BC         & \\
1C14 & F1       &    & POP AF         & \\
1C15 & 23       &    & INC HL         & \\
1C16 & 18 EE    &    & JR 1E          & \\
1C18 & 42 4F \dots & TB & DEFM "BONJOUR DE X07" & message \\
1C26 & 00       &    & DEFB \$00      & fin du message \\
\end{asm}

\paragraph{Chargeur BASIC.}
\begin{verbatim}
10 CLS:PRINT"un instant !"
20 FOR I=&H1C00 TO &H1C26
30 READ A$:POKE I,VAL("&H"+A$)
40 NEXT I
60 DATA CD,B7,CF,21,18,1C,7E,FE,00,C8,F5,C5,D5,E5,CD,D6,CE,E1,D1,C1,F1,23
70 DATA 18,EE,42,4F,4E,4A,4F,55,52,20,44,45,20,58,30,37,00
\end{verbatim}
Le chargeur ne lance pas la routine : taper \texttt{EXEC \&H1C00}, la X-710
étant branchée et allumée.

%======================================================================
\section{Bruitages \pg{131-132}}

\paragraph{But.} Produire avec le buzzer des sons impossibles à obtenir
avec \texttt{BEEP}, en programmant directement les ports d'E/S. Le livre
rappelle un article de la gazette n\textsuperscript{o}\,1 du Club C7 qui
produisait « des sons assez inattendus ».

\paragraph{Principe.} Le programme utilise trois ports du circuit d'E/S
(voir §\,7.5 du livre et la section~\ref{sec:ports}) :
\begin{itemize}
  \item \textbf{port \hex{F4}} : on y écrit \hex{4E}, soit BRGST\,=\,1
        (le générateur de bauds compte), MD1\,=\,MD0\,=\,1 (sortie sur le
        buzzer) et BZON\,=\,1 (buzzer en marche) ;
  \item \textbf{ports \hex{F2} et \hex{F3}} : ils fixent la fréquence des
        impulsions. C'est un compteur de 12 bits : 8 bits dans \hex{F2},
        4 bits dans \hex{F3}.
\end{itemize}
La routine fait varier \hex{F2} de \hex{FF} à 1 pour chaque valeur de
\hex{F3} (d'abord 0, puis de \hex{0E} à 1). Une petite temporisation (TEMPO)
rend les sons audibles. Pour arrêter le buzzer, on écrit \hex{00} dans
\hex{F4}.

\fiche{\hex{1C00}--\hex{1C2B}}{44 octets}{\texttt{EXEC \&H1C00}}

\begin{asm}
1C00 & 3E 00    &    & LD A,\$00      & \\
1C02 & 06 0F    &    & LD B,\$0F      & valeur de départ de F3 \\
1C04 & D3 F3    &    & OUT (\$F3),A   & \\
1C06 & 3E 4E    &    & LD A,\$4E      & BRGST, MD1, MD0, BZON \\
1C08 & D3 F4    &    & OUT (\$F4),A   & buzzer en marche \\
1C0A & 3E FF    & 2E & LD A,\$FF      & \\
1C0C & D3 F2    & 1E & OUT (\$F2),A   & poids faible de la fréquence \\
1C0E & 3D       &    & DEC A          & \\
1C0F & 28 05    &    & JR Z,F3        & \\
1C11 & CD 24 1C &    & CALL TE        & temporisation \\
1C14 & 18 F6    &    & JR 1E          & \\
1C16 & 78       & F3 & LD A,B         & \\
1C17 & 3D       &    & DEC A          & \\
1C18 & 28 05    &    & JR Z,FI        & \\
1C1A & D3 F3    &    & OUT (\$F3),A   & poids fort de la fréquence \\
1C1C & 47       &    & LD B,A         & \\
1C1D & 18 EB    &    & JR 2E          & \\
\multicolumn{5}{@{}l}{\small\itshape Arrêt} \\
1C1F & 3E 00    & FI & LD A,\$00      & \\
1C21 & D3 F4    &    & OUT (\$F4),A   & buzzer arrêté \\
1C23 & C9       &    & RET            & \\
\multicolumn{5}{@{}l}{\small\itshape Temporisation} \\
1C24 & F5       & TE & PUSH AF        & \\
1C25 & 3E 80    &    & LD A,\$80      & \\
1C27 & 3D       & LO & DEC A          & \\
1C28 & 20 FD    &    & JR NZ,LO       & \\
1C2A & F1       &    & POP AF         & \\
1C2B & C9       &    & RET            & \\
\end{asm}

\paragraph{Chargeur BASIC.}
\begin{verbatim}
10 PRINT"un instant !"
20 FOR I=&H1C00 TO &H1C2B
30 READ A$:POKE I,VAL("&H"+A$)
40 NEXT
50 EXEC&H1C00
60 DATA 3E,00,06,0F,D3,F3,3E,4E,D3,F4,3E,FF,D3,F2,3D,28,05,CD,24,1C
70 DATA 18,F6,78,3D,28,05,D3,F3,47,18,EB,3E,00,D3,F4,C9,F5,3E,80,3D,20,FD,F1,C9
\end{verbatim}
On obtient d'autres sons en changeant les valeurs de départ (\hex{1C01},
\hex{1C03}, \hex{1C0B}) ou la durée de la temporisation (\hex{1C26}).

\nota{Le livre parle des « ports de sortie du NSC800 » : ce sont en fait les
ports du circuit d'E/S HD61L202F. Il écrit aussi « 8 octets dans le port
F2 et 4 autres dans le port F3 » : il s'agit de bits.}

%======================================================================
\section{Redéfinition d'une touche \pg{133-134}}

\paragraph{But.} Associer un message à une touche du clavier. Ici, la
touche \textasciitilde{} (tilde) affiche « X07 ». La même méthode permet de
redéfinir tout le clavier, par exemple en AZERTY.

\paragraph{Principe.} Tout caractère affiché passe par la sortie console,
appelée par le saut placé en \hex{009F} dans la zone système (\texttt{JP
C1BE}, voir §\,8.3 du livre). La routine se compose de deux parties :
\begin{enumerate}
  \item l'installation remplace l'adresse de ce saut (octets \hex{00A0} et
        \hex{00A1}) par celle de la routine personnelle ;
  \item la routine personnelle compare le caractère reçu dans \reg{A} au
        code de la touche redéfinie (\hex{7E}). S'ils sont égaux, elle
        affiche le message par \hex{D5B0} et revient au mode « attente
        curseur » ; sinon, elle passe la main à la routine d'origine
        \hex{C1BE}, qui affiche le caractère normalement.
\end{enumerate}
L'organigramme n'a pas de fin : une fois installée, la routine est appelée
à chaque caractère affiché et décide elle-même si elle doit agir.

\fiche{\hex{1C00}--\hex{1C1B} (code) ; \hex{1C1C}--\hex{1C20} (message)}%
      {28 octets de code + 5 octets de message}{lancé par le chargeur}

\begin{asm}
\multicolumn{5}{@{}l}{\small\itshape Changement du saut (hook)} \\
1C00 & 21 07 1C &    & LD HL,DB       & adresse de la routine personnelle \\
1C03 & 22 A0 00 &    & LD (\$00A0),HL & remplace l'adresse du JP en \hex{009F} \\
1C06 & C9       &    & RET            & \\
\multicolumn{5}{@{}l}{\small\itshape Routine personnelle} \\
1C07 & F5       & DB & PUSH AF        & \\
1C08 & FE 7E    &    & CP "\textasciitilde" & touche redéfinie ? \\
1C0A & 28 04    &    & JR Z,AF        & \\
1C0C & F1       &    & POP AF         & \\
1C0D & C3 BE C1 &    & JP \$C1BE      & affichage normal du caractère \\
1C10 & E5       & AF & PUSH HL        & \\
1C11 & 21 1C 1C &    & LD HL,ME       & \\
1C14 & CD B0 D5 &    & CALL \$D5B0    & affichage du message \\
1C17 & E1       &    & POP HL         & \\
1C18 & F1       &    & POP AF         & \\
1C19 & C3 3D F2 &    & JP \$F23D      & attente curseur \\
1C1C & 20 58 30 37 & ME & DEFM " X07" & message \\
1C20 & 00       &    & DEFB \$00      & fin du message \\
\end{asm}

\paragraph{Chargeur BASIC.}
\begin{verbatim}
10 PRINT"un instant"
20 FOR I=&H1C00 TO &H1C20
30 READ A$:POKE I,VAL("&H"+A$)
40 NEXT
50 EXEC&H1C00:PRINT"APPUYEZ SUR ~"
60 DATA 21,07,1C,22,A0,00,C9,F5,FE,7E,28,04,F1,C3,BE,C1,E5,21,1C,1C,CD,B0,D5
70 DATA E1,F1,C3,3D,F2,20,58,30,37,00
\end{verbatim}

\paragraph{Remarques.}
\begin{itemize}
  \item Comme la routine intercepte la sortie console, le message apparaît
        aussi quand un programme affiche le caractère \textasciitilde, par
        exemple \texttt{PRINT "\textasciitilde"}. C'est d'ailleurs ce que
        fait la ligne 50 du chargeur.
  \item Pour revenir au fonctionnement normal, il faut remettre l'adresse
        d'origine : \texttt{POKE \&HA0,\&HBE:POKE \&HA1,\&HC1}.
\end{itemize}
\nota{Le texte indique que la routine d'origine est en \hex{CEB1}, mais le
code saute en \hex{C1BE}, l'adresse du saut d'origine en \hex{009F}. Le
livre annonce aussi 32 octets (\hex{1C00}--\hex{1C1F}) : le programme
occupe en réalité 33 octets, jusqu'à \hex{1C20}.}

%======================================================================
\section{Conclusion du livre \pg{135}}

La 3\textsuperscript{e} partie se termine sur un court récit daté
d'octobre 1985 : CANDORE, qui s'était endormi devant son X-07, se réveille
et voit l'écran clignoter, le buzzer émettre des sons inconnus du BASIC,
l'imprimante tracer à grande vitesse et le téléviseur afficher
« EUREKA ». Grâce au « vieux sage ASSEMBLUS » et à l'ouvrage \emph{le
Canon X-07 et l'assembleur}, il a découvert tous les secrets de sa
machine. Le texte est signé A.~Tonic.

Viennent ensuite les annexes, qui ne sont pas reprises ici : liste des codes
Z80 (annexe 1), assembleurs et désassembleurs disponibles (annexe 2),
présentation du Club C7 (annexe 3) et bibliographie (annexe 4).

\end{document}
